%% 第26回編集委員会（物理第4巻）：組 電W1
%% p4-01.tex  電-1　電気の力と場：クーロンの法則と電場（★1〜4・3段にしない）
%% 中身は計画書 tools/phys4-plan.md の B「電-1」。公式の持ち主は A（A10・A20・A21・A60）。
%% 数値はすべて Python で計算した：checks/物理第4巻/電-1/calc_p4_01.py（出力 calc_p4_01_out.txt）、
%%   fieldlines_p4_01.py（⑤の電気力線の座標）、map_p4_01.py（⑩の表と矢印の地図）。
%% 前の巻との約束（計画書 A-6）：
%%   ・P1（力-2 の表「電気・磁石の力は『電磁気学』で」）→ ②（クーロンの法則。磁石の力は電-6）
%%   ・P9（第1巻の巻末：電場 E・電荷 Q は電-1）→ ①④の記号カード
%%   ・P21・P41（第2・3巻の巻末：電荷 Q・電場 E は電-1）→ 同じカード
%% 数学との約束（計画書 C-6-2）：8-5⑥ の表（1 nC：899・225・99.9・36.0・8.99、k＝8.99×10⁹）→ ④・例題7（数を変えずに使った）
%%   8-11⑦ の平面 σ/(2ε0)・線 λ/(2πε0 r) は、ここでは重ね合わせの積分（8-6）で導き、ガウスの法則による導き方は電-2 ④で「楽な道」として見せる。
%% 記号（この章）：電荷 q（試しの・小さな電荷）・Q（場をつくる電荷）、電場 \vec E（大きさ E）、クーロンの定数 k（ばね定数ではない）、
%%   真空の誘電率 ε0、電気素量 e、線電荷密度 λ、面電荷密度 σ、双極子モーメント \vec p（大きさ p＝qd。運動量は出さない）、
%%   距離 r、双極子の電荷の間かく d、円環の半径 a、円板の半径 R、軸の上の位置 z、質量 m、加速度 a（円環の半径と同じ節で使わない）。
%%   面積は S（Y-4。この章では使わない）。
%% 出典（11月5日に確かめた）
%%   ・CODATA 2022（NIST physics.nist.gov/cuu を curl で読んだ）：ε0＝8.854 187 8188×10⁻¹² F/m、e＝1.602 176 634×10⁻¹⁹ C（定義値）、
%%       mₑ＝9.109 383 7139×10⁻³¹ kg、mₚ＝1.672 621 925 95×10⁻²⁷ kg、ボーア半径 5.291 772 105 44×10⁻¹¹ m、G＝6.674 30×10⁻¹¹。
%%       1/(4πε0)＝8.987 55×10⁹（本文は Y-6 の 8.99×10⁹）。μ0 が測る値になった 2019 年の SI の改定で、ε0＝1/(μ0c²) も測る値になった。
%%   ・クーロンのねじればかりの実験（1785年、パリの王立科学アカデミーの『第1論文』）：National MagLab "Torsion Balance - 1785"
%%       （nationalmaglab.org の博物館のページ）と、W. Whewell, History of the Inductive Sciences, Vol. II（1858）の該当ページ
%%       （19世紀の科学書の複写 geology.19thcenturyscience.org）。WebSearch の抜粋で2つを照らし合わせた。
%%   ・こすったときに何が移るのか（電子・イオン・物のかけら）は、長い研究のあとも、はっきりしていない：
%%       D. J. Lacks, T. Shinbrot, "Long-standing and unresolved issues in triboelectric charging", Nat. Rev. Chem. 3, 465–476 (2019)、
%%       DOI 10.1038/s41570-019-0115-1（書誌と要旨を WebSearch の抜粋で確かめた。本文は有料で読んでいない）、
%%       J. Haeberle ほか "Double origin of stochastic granular tribocharging", Soft Matter (2018)、arXiv:1803.05280（要旨を arxiv.org で読んだ。
%%       帯電はばらつきが大きく、同じ条件でも両方の符号が出る）。
%%   ・雷1回で地面へ運ばれる電荷：CIGRE の雷のパラメータ（TB 549）の中央値（負の雷の1回めの放電 約5 C、続く放電 約1.4 C、正の雷 約80 C）を
%%       引く資料の WebSearch の抜粋と、Rakov & Uman の教科書の記述（数 C の電荷が地面へ）で照らし合わせ、本文は「数 C から数十 C」とした。
%%   ・人がパチッと感じる体の電圧 約3 kV（3000 V）：作業服の帯電の研究（KAKEN 15K00750 の報告）の WebSearch の抜粋と、
%%       大阪産業技術研究所の技術資料（人体の静電容量 約100 pF・3 kV の計算例）の抜粋で照らし合わせた。どちらも作業環境から本文には届かなかった。
%%   ・水の比誘電率 約80（20 ℃で80.2）：計画書 C-3〔確〕（IAPWS R8-97 の式）。
%%   ・水分子の双極子モーメント 約6.2×10⁻³⁰ C·m（1.85 デバイ）は本文で使わない（参考の数として calc に置いた）。
%% 〔要確認〕なし（髪の毛・インクジェットの数は、本文で「〜とします」と書いた仮の数）。
\usetikzlibrary{decorations.markings}
\makeatletter
% 点電荷（丸の中に＋・−。色だけで分けない）  \sxbplus{(中心)}{半径}  \sxbminus{(中心)}{半径}
\newcommand\sxbplus[2]{\filldraw[fill=bkRed!75!black, draw=bkInk, line width=0.5pt] #1 circle (#2); \node[font=\sffamily\bfseries\footnotesize, text=white, inner sep=0pt] at #1 {$+$};}
\newcommand\sxbminus[2]{\filldraw[fill=bkBlue!75!black, draw=bkInk, line width=0.5pt] #1 circle (#2); \node[font=\sffamily\bfseries\footnotesize, text=white, inner sep=0pt] at #1 {$-$};}
% 小さな＋・−の印（原子・物の中の電荷）  \sxbp{(位置)}  \sxbm{(位置)}
\newcommand\sxbp[1]{\node[font=\sffamily\bfseries\footnotesize, text=bkRed!75!black, inner sep=0pt] at #1 {$+$};}
\newcommand\sxbm[1]{\node[font=\sffamily\bfseries\footnotesize, text=bkBlue!75!black, inner sep=0pt] at #1 {$-$};}
% 図の中の字・電気力線・電場の矢印・力の矢印
\tikzset{sxb lab/.style={font=\sffamily\footnotesize, text=bkInk},
  sxb line/.style={bkPurple!80!black, line width=0.7pt},
  sxb tip/.style={bkPurple!80!black, line width=0.7pt, -{Stealth[length=1.9mm, width=1.5mm]}},
  sxb arr/.style={bkNavy, line width=0.7pt, -{Stealth[length=1.5mm, width=1.2mm]}},
  sxb f/.style={line width=1.2pt, -{Stealth[length=2.6mm, width=2.0mm]}},
  sxb dim/.style={bkGray, line width=0.5pt, {Stealth[length=1.6mm]}-{Stealth[length=1.6mm]}}}
\makeatother

\phchapter{1}{電気の力と場：\\クーロンの法則と電場}{下じきでこすった髪の毛は、なぜ立ち上がる？}{\bklevelrange{1}{4}}{電気の力と場：クーロンの法則と電場}
\bkmichishirube{4-10,5-7,5-17,8-5,8-6,力-1,力-2}{電-2,電-3,電-4,電-6,電-15}

\begin{goal}
\begin{itemize}[label={\color{bkTheme}●}, leftmargin=1.6\zw, itemsep=0pt, topsep=1pt]
\item \bklevel{1}\enspace こすると物が電気を帯びるわけを、電子が移ることで説明できる。
\item \bklevel{3}\enspace クーロンの法則と重ね合わせで、いくつかの電荷のあいだの力と、電場を求められる。
\item \bklevel{3}\enspace 一様な電場の中の電荷の動きを、力学の放物運動と同じ考えで求められる。
\item \bklevel{4}\enspace 線・円環・平面に広がった電荷の電場を積分で求め、電気双極子の電場が距離の3乗に反比例して弱まることを示せる。
\end{itemize}
\end{goal}

\phkoushiki{%
\phk{クーロンの法則 $F=k\,\dfrac{q_1q_2}{r^2}$、$k=\dfrac{1}{4\pi\varepsilon_0}$}{高校物理}{ph:電-1:coulomb}
\phk{電場 $\vec E=\dfrac{\vec F}{q}$、点電荷の電場 $E=k\,\dfrac{Q}{r^2}$、重ね合わせ}{高校物理}{ph:電-1:tenkaba}
\phkx{電気力線の本数 $4\pi kQ$}{高校物理}{電-2 で導く}
}

\begin{hatena}
\noindent\begin{minipage}[c]{0.55\linewidth}\raggedright
乾いた冬の日に、プラスチックの下じきで頭をこすってから、下じきをゆっくり持ち上げてみましょう。髪の毛が、下じきに引っぱられるように、ふわっと立ち上がります。下じきは髪に触れていません。それなのに、髪の毛を持ち上げているのは、どんな力でしょう。

髪の毛1本は軽いとはいえ、地球が下へ引く重力を受けています。それに負けないほどの力が、離れた下じきから届いているのでしょうか。
\end{minipage}\hfill
\begin{minipage}[c]{0.42\linewidth}\centering
\begin{tikzpicture}[x=0.88mm, y=0.88mm]
  % 下じき
  \filldraw[fill=bkBlue!18, draw=bkBlue!70!black, line width=0.7pt, rotate around={-8:(24,44)}] (2,42) rectangle (46,46);
  \foreach \xx in {8,14,...,40} {\sxbm{({\xx},{44.2-0.14*(\xx-24)})}}
  % 頭
  \filldraw[fill=bkOrange!25, draw=bkOrange!70!black, line width=0.7pt] (24,6) circle (11);
  \foreach \a in {35,50,...,145} {\draw[bkInk!80, line width=0.6pt] ($(24,6)+(\a:11)$) .. controls ($(24,6)+(\a:16)$) and ($(24,6)+({90+0.55*(\a-90)}:22)$) .. ($(24,6)+({90+0.45*(\a-90)}:{28-0.04*abs(\a-90)})$);}
  \sxbp{(16.6,34.6)}\sxbp{(31.4,34.6)}\sxbp{(24,36.6)}
  \node[sxb lab, anchor=west] at (47,38) {下じき};
  \node[sxb lab, anchor=west] at (36,16) {髪の毛};
\end{tikzpicture}
\end{minipage}
\end{hatena}

\noindent 電気というと、コンセントや電池を思いうかべるかもしれません。でも、電気のはじまりは、こすった物どうしが引き合う「静電気」の不思議でした。2500年以上前のギリシャの人々も、こすったこはく（木のやにの化石）がちりを引きつけることを知っていました。英語の electricity（電気）は、ギリシャ語の「こはく」から来ています。この章では、離れた物のあいだにはたらく電気の力を、数で表す方法を学びます。そして、電気の力を「電荷の周りの空間がもつ性質」として表す\textbf{電場}という考えを作ります。電場は、この巻の最後に出てくる電波・光（電-15）までつながる、電磁気学のいちばんの土台です。

\begin{yattemiyou}[静電気で、引きつける・はね返す]
\textbf{用意するもの}：下じき、ティッシュペーパー、プラスチックのストロー2本、糸、水道\par\vspace{2pt}
\noindent\begin{minipage}[c]{0.62\linewidth}
\begin{enumerate}[label={\bkcircled[bkGreen]{\arabic*}}, leftmargin=2.4\zw, itemsep=1pt, topsep=0pt]
\item 下じきを髪の毛で10回ほどこすり、ちぎったティッシュや、細く出した水道の水に近づける。
\item ストロー2本を、それぞれティッシュでこする。1本を糸でつるし、もう1本を近づける。
\item こすったティッシュを、つるしたストローに近づける。
\item 乾いた日と、しめった日でくらべる。
\end{enumerate}
\end{minipage}\hfill
\begin{minipage}[c]{0.35\linewidth}\centering
\begin{tikzpicture}[x=1mm, y=1mm]
  \fill[bkGray!35] (0,26) rectangle (24,28);
  \draw[bkInk, line width=0.6pt] (0,26) -- (24,26);
  \draw[bkInk, line width=0.5pt] (11,26) -- (11,15);
  \draw[bkOrange!80!black, line width=2.2pt] (0,15) -- (22,15);
  \foreach \xx in {3,8,...,20} {\sxbm{(\xx,12.4)}}
  \draw[bkOrange!80!black, line width=2.2pt] (30,15) -- (44,15);
  \foreach \xx in {33,38,43} {\sxbm{(\xx,12.4)}}
  \draw[bkRed!80!black, line width=1pt, -{Stealth[length=2mm]}] (22,17.5) arc (-20:30:7);
  \draw[bkInk, line width=0.8pt, -{Stealth[length=2mm]}] (40,19) -- (32,19);
  \node[sxb lab, anchor=north, align=center] at (11,9) {つるした\\ストロー};
  \node[sxb lab, anchor=north, align=center] at (38,9) {近づける\\ストロー};
\end{tikzpicture}
\end{minipage}
\end{yattemiyou}

\noindent やってみようの答え：①ティッシュの小さな切れはしは、下じきに飛びついてきます。水の細い流れは、下じきのほうへ曲がります。②こすったストローどうしは、近づけると、つるしたほうが逃げるように回ります（\textbf{はね返す}）。③こすったティッシュを近づけると、つるしたストローは寄ってきます（\textbf{引きつける}）。④しめった日は、どれも弱くなります。

同じ物を同じようにこすったストロー2本は「同じ種類の電気」をもっていて、はね返し合います。ストローとこすったティッシュは「違う種類の電気」をもっていて、引き合います。引き合う・はね返す、の2通りがあることが、重力（いつも引き合うだけ）との大きな違いです。この章の①と②で、そのしくみと力の大きさを調べます。①の水が曲がるわけは⑥で説明します。

\begin{omoidasou}
\bkrecall{4-10}{ベクトルのたし算は、矢印をつなぐ（成分ごとにたす）。$\vec a=(a_x,\ a_y)$ の大きさは $\sqrt{a_x^2+a_y^2}$。}
\bkrecall{8-5}{場所ごとに矢印が1本決まるものをベクトル場という。矢印の向きをつないだ線が力線。点電荷の周りの電場の大きさは $E=kQ/r^2$、$k=8.99\times10^9$\,N·m²/C²（8-5⑥）。}
\bkrecall{力-2}{運動方程式 $m\vec a=\vec F$。万有引力 $F=GMm/r^2$（$G\approx\qty{6.67e-11}{N.m^2/kg^2}$）。}
\bkrecall{力-1}{横の速さが一定で、たてに一定の加速度 $a$ を受ける物は、放物線をえがく。$t$ 秒でたてに $\tfrac12at^2$ ずれる。}
\bkrecall{8-6}{広がった物の量は、小さく分けて、1つずつかけてたす（積分）。}
\bkrecall{5-17}{$x$ が1よりずっと小さいとき、$(1+x)^n\approx1+nx$。}
\end{omoidasou}

\noindent この章の記号：電荷（電気の量）を $q$・$Q$ と書きます。力を受ける小さな電荷を $q$、周りに電場をつくる電荷を $Q$ と書き分けることが多いですが、どちらも同じ「電荷」です。\textbf{クーロンの定数} $k$ は、第1巻のばね定数 $k$ とは別の量です（この巻の電磁気の章では、ばね定数を使いません）。

%% ======================== 物理の道具 ========================
\bksection{doguu}{物理の道具}

\bksubsection{1}{静電気と電荷}
\label{ph:電-1:sec1}
どんな物も、とても小さな\textbf{原子}が集まってできています。原子のまん中には、\textbf{正}（プラス）の電気をもつ原子核があり、その周りを、\textbf{負}（マイナス）の電気をもつ\textbf{電子}がとり巻いています。ふだんの物は、正の電気と負の電気がちょうど同じ量だけあって、打ち消し合っています。外から見ると、電気をもっていないのと同じです。

\begin{bkfigure}
\begin{tikzpicture}[x=0.94mm, y=0.94mm]
  % こする前
  \begin{scope}
    \filldraw[fill=bkOrange!12, draw=bkOrange!70!black, line width=0.6pt] (0,16) rectangle (34,26);
    \filldraw[fill=bkBlue!10, draw=bkBlue!70!black, line width=0.6pt] (0,2) rectangle (34,12);
    \foreach \xx in {4,10,...,28} {\sxbp{(\xx,23)}\sxbm{({\xx+2.5},19)}\sxbp{(\xx,9)}\sxbm{({\xx+2.5},5)}}
    \node[sxb lab, anchor=west] at (35,21) {髪の毛};
    \node[sxb lab, anchor=west] at (35,7) {下じき};
    \node[sxb lab, anchor=north] at (17,0) {こする前：＋と－が同じ数};
  \end{scope}
  \draw[bkInk, line width=1pt, -{Stealth[length=3mm]}] (50,14) -- (62,14);
  \node[sxb lab, anchor=south] at (56,15) {こする};
  % こすった後
  \begin{scope}[shift={(70,0)}]
    \filldraw[fill=bkOrange!12, draw=bkOrange!70!black, line width=0.6pt] (0,20) rectangle (34,30);
    \filldraw[fill=bkBlue!10, draw=bkBlue!70!black, line width=0.6pt] (0,0) rectangle (34,10);
    \foreach \xx in {4,10,...,28} {\sxbp{(\xx,27)}\sxbp{(\xx,8)}\sxbm{({\xx+2.5},2)}}
    \foreach \xx in {6.5,18.5,30.5} {\sxbm{(\xx,23)}}
    \foreach \xx in {12.5,24.5} {\sxbm{(\xx,5)}}
    \draw[bkBlue!70!black, line width=0.7pt, -{Stealth[length=2mm]}] (12.5,22) -- (12.5,10.6);
    \draw[bkBlue!70!black, line width=0.7pt, -{Stealth[length=2mm]}] (24.5,22) -- (24.5,10.6);
    \node[sxb lab, anchor=west, align=left] at (35,25) {髪の毛：電子が減って\\正に帯電};
    \node[sxb lab, anchor=west, align=left] at (35,5) {下じき：電子が増えて\\負に帯電};
    \node[sxb lab, anchor=north] at (17,-2) {電子が移る};
  \end{scope}
\end{tikzpicture}
\bkfigcaption{こすると、髪の毛の電子の一部が下じきへ移る。電気が生まれたのではなく、移っただけ。髪の毛の＋が多い分と、下じきの－が多い分は、同じ量。}\label{ph:電-1:fig-kosuru}
\end{bkfigure}

\noindent 2つの物をこすり合わせると、ふれ合った所で、一方の物の電子の一部が、もう一方の物へ移ります（図\ref{ph:電-1:fig-kosuru}）。電子をもらった物は、負の電気が多くなり、\textbf{負に帯電}します。電子を取られた物は、正の電気が多く残り、\textbf{正に帯電}します。\bkidx{たいでん}{帯電}髪の毛とプラスチックの下じきでは、髪の毛から下じきへ電子が移り、髪の毛は正に、下じきは負に帯電します。このように、物にたまって動かない電気を\textbf{静電気}といいます。\bkidx{せいでんき}{静電気}

\phboxhead{電荷と、その単位}
物がもっている電気の量を\textbf{電荷}（または電気量）といい、正の電荷は正の数、負の電荷は負の数で表します。\bkidx{でんか}{電荷}単位は\textbf{クーロン}〔C〕です。1\,C は、\qty{1}{A}（アンペア）の電流が1秒間流れたときに運ばれる電荷です（電流は電-4）。電子1個の電荷は $-e$ で、
\pheq{e=\qty{1.602176634e-19}{C}\approx\qty{1.60e-19}{C}}{C}
です。$e$ を\textbf{電気素量}といいます。\bkidx{でんきそりょう}{電気素量}陽子（原子核の中の正の粒）1個の電荷は、ちょうど $+e$ です。2019年からは、この値を「ちょうどこの数」と決めて、それをもとに A や C の大きさが決められています。1\,C は $1\div(1.602\times10^{-19})\approx6.24\times10^{18}$ 個、つまり624京個の電子の電荷の大きさです。

\phboxhead{電荷は、なくならない}
こすり合わせても、電子が一方からもう一方へ移るだけなので、2つの物の電荷を合わせた量は、こする前と同じ0のままです。電荷は、勝手に生まれたり消えたりしません。これを\textbf{電荷の保存}といいます。\bkidx{でんかのほぞん}{電荷の保存}同じ形の2つの金属球の一方が $+3$\,nC、他方が $-1$\,nC を帯びていて、2つを触れさせてから離すと、合計の $+2$\,nC が半分ずつ分かれて、どちらも $+1$\,nC になります（形が同じなので、ちょうど半分ずつ。例題1のじぶんでとこう）。1\,nC（ナノクーロン）は $10^{-9}$\,C です。

\phboxhead{電気を通す物・通さない物}
金属の中には、原子から離れて、物の中を自由に動ける電子（\textbf{自由電子}）がたくさんあります。だから、金属のはしに電荷をあたえると、電荷はすぐに金属全体へ広がります。このように電気をよく通す物を\textbf{導体}、ガラス・ゴム・プラスチックのように電子がほとんど動けない物を\textbf{絶縁体}（不導体）といいます。\bkidx{どうたい}{導体}\bkidx{ぜつえんたい}{絶縁体}こすって静電気をためやすいのは、電荷がその場に留まる絶縁体です。手に持った金属の棒をこすっても、電荷は手と体を通って地面へにげてしまいます（導体の性質は電-3、電流は電-4）。

\begin{bkfigure}
\begin{tikzpicture}[x=1mm, y=1mm]
  \shade[left color=bkRed!30, right color=bkBlue!30] (0,0) rectangle (124,8);
  \draw[bkInk, line width=0.5pt] (0,0) rectangle (124,8);
  \node[sxb lab, anchor=west] at (1,4) {＋になりやすい};
  \node[sxb lab, anchor=east] at (123,4) {－になりやすい};
  \foreach \xx/\nm/\yy in {12/髪の毛・羊毛/11, 36/ナイロン/11, 54/絹/11, 72/紙・木綿/11, 94/ポリエチレン/16, 117/塩化ビニル/11} {\draw[bkInk, line width=0.5pt] (\xx,8) -- (\xx,\yy); \node[sxb lab, anchor=south] at (\xx,\yy) {\nm};}
\end{tikzpicture}
\bkfigcaption{帯電列の例。2つの物をこすると、左にある物が正、右にある物が負に帯電しやすい。順番は、資料や、こするときの条件（湿度・よごれ・こすり方）で入れかわることがある。}\label{ph:電-1:fig-retsu}
\end{bkfigure}

\noindent どちらの物が正になり、どちらが負になるかは、物の組み合わせで決まります。こすったときに正になりやすい順に並べたものを\textbf{帯電列}といいます（図\ref{ph:電-1:fig-retsu}）。\bkidx{たいでんれつ}{帯電列}同じ物どうし（ストロー2本）は同じ符号に帯電し、帯電列で離れた物どうしほど、強く帯電する傾向があります。しめった日に静電気が起きにくいのは、物の表面についた薄い水の膜を通って、電荷が少しずつにげていくからです。

\begin{phmada}[こすると、何が移るのか？]
こすり合わせると帯電することは、2500年以上前から知られています。ところが、\textbf{何が}移っているのか（電子なのか、表面にあるイオンなのか、こすれて取れた物の小さなかけらなのか）は、今もはっきりしていません。組み合わせによって、移るものが違うとも考えられています。帯電列の順番が実験によって入れかわることや、同じ条件でこすっても帯電の量が大きくばらつき、ときには符号まで逆になることも、まだうまく説明できていません。粉の工場の事故（粉じん爆発）を防いだり、静電気で発電する装置を作ったりするために、物理・化学・機械の研究者がいっしょに研究を続けています。この章では、「電子が移る」という、いちばん簡単な考え方で話を進めます。
\end{phmada}

%% ---------------------------------------------------------------
\bksubsection{2}{クーロンの法則：電気の力の大きさ}
\label{ph:電-1:sec2}
帯電した物どうしの力は、どれくらいの大きさで、距離によってどう変わるのでしょう。フランスの物理学者クーロンは、1785年、細い針金のねじれで小さな力をはかる\textbf{ねじればかり}を作って、この問題を調べました（図\ref{ph:電-1:fig-nejiri}）。

\begin{bkfigure}
\begin{tikzpicture}[x=1mm, y=1mm]
  % ガラスの筒
  \draw[bkGray!80!black, line width=0.6pt] (10,0) -- (10,24) (70,0) -- (70,24) (10,0) -- (70,0);
  \draw[bkGray!80!black, line width=0.6pt] (10,24) -- (70,24);
  \draw[bkGray!80!black, line width=0.6pt] (36,24) -- (36,44) (44,24) -- (44,44);
  \fill[bkGray!40] (34,44) rectangle (46,47);
  \draw[bkInk, line width=0.5pt] (34,44) rectangle (46,47);
  \node[sxb lab, anchor=west] at (47,45.5) {ねじりの目もり};
  % 針金
  \draw[bkInk, line width=0.4pt] (40,44) -- (40,14);
  % 棒と球
  \draw[bkInk, line width=0.9pt] (18,14) -- (62,14);
  \filldraw[fill=bkGray!30, draw=bkInk, line width=0.5pt] (18,14) circle (1.4);
  \sxbplus{(58,14)}{2.2}
  \sxbplus{(48,7)}{2.2}
  \draw[bkInk, line width=0.5pt] (48,9.2) -- (48,24);
  \draw[bkRed!80!black, sxb f] (59,15.5) -- (65,19);
  \node[sxb lab, anchor=west] at (65.5,20) {はね返す力};
  \node[sxb lab, anchor=west, align=left] at (73,7) {針金がねじれて\\もどろうとする力と\\つり合う所で止まる};
  \node[sxb lab, anchor=east] at (8.5,6) {つり合いのおもり};
  \draw[bkInk!60, line width=0.4pt] (8.8,6.5) -- (17,13);
  \node[sxb lab, anchor=north] at (48,4.6) {固定した球};
  \node[sxb lab, anchor=east] at (35,32) {細い針金};
\end{tikzpicture}
\bkfigcaption{クーロンのねじればかり（しくみの図）。同じ符号に帯電した2つの小さな球がはね返し合うと、棒が回って針金がねじれる。ねじれの角から力の大きさが、球の位置から距離がわかる。}\label{ph:電-1:fig-nejiri}
\end{bkfigure}

\noindent クーロンは、2つの球の距離を半分にすると、はね返す力がおよそ4倍になることを見つけました。また、帯電した金属球に、同じ大きさの帯電していない金属球を触れさせると、電荷はちょうど半分ずつに分かれます（①の電荷の保存）。こうして電荷を半分にすると、力も半分になります。このような実験から、次の法則がわかりました。

\begin{bkfigure}
\begin{tikzpicture}[x=1mm, y=1mm]
  \sxbplus{(10,20)}{2.4}\sxbplus{(40,20)}{2.4}
  \draw[bkRed!80!black, sxb f] (7.4,20) -- (-2,20);
  \draw[bkRed!80!black, sxb f] (42.6,20) -- (52,20);
  \draw[sxb dim] (10,9) -- (40,9);
  \node[sxb lab, anchor=north] at (25,8.5) {$r$};
  \node[sxb lab, anchor=south] at (25,23) {同じ符号：はね返す};
  \sxbplus{(80,20)}{2.4}\sxbminus{(110,20)}{2.4}
  \draw[bkRed!80!black, sxb f] (82.6,20) -- (92,20);
  \draw[bkRed!80!black, sxb f] (107.4,20) -- (98,20);
  \draw[sxb dim] (80,9) -- (110,9);
  \node[sxb lab, anchor=north] at (95,8.5) {$r$};
  \node[sxb lab, anchor=south] at (95,23) {違う符号：引き合う};
  \node[sxb lab, anchor=north] at (10,17) {$q_1$};
  \node[sxb lab, anchor=north] at (40,17) {$q_2$};
  \node[sxb lab, anchor=north] at (80,17) {$q_1$};
  \node[sxb lab, anchor=north] at (110,17) {$q_2$};
\end{tikzpicture}
\bkfigcaption{2つの点電荷のあいだの力。向きは2つを結ぶ線の上。2つの電荷は、同じ大きさの力を、反対向きに受ける（作用・反作用：力-2）。}\label{ph:電-1:fig-chikara}
\end{bkfigure}

\noindent 距離 $r$ 離れた2つの小さな電荷 $q_1$・$q_2$ のあいだにはたらく力の大きさ $F$ は、
\pheq{F=k\,\frac{|q_1q_2|}{r^2}}{N ＝ $\dfrac{\text{N·m}^2}{\text{C}^2}$×$\dfrac{\text{C}\times\text{C}}{\text{m}^2}$}
\label{ph:電-1:coulomb}で、同じ符号ならはね返し合い、違う符号なら引き合います。これを\textbf{クーロンの法則}といいます。\bkidx{くーろんのほうそく}{クーロンの法則}比例の定数 $k$ は、
\pheq{k=\frac{1}{4\pi\varepsilon_0}\approx\qty{8.99e9}{N.m^2/C^2},\qquad \varepsilon_0\approx\qty{8.85e-12}{F/m}}{$\dfrac{1}{\text{F/m}}$ ＝ N·m²/C²}
です。$k$ を\textbf{クーロンの定数}、$\varepsilon_0$（イプシロン・ゼロ）を\textbf{真空の誘電率}といいます。\bkidx{くーろんのていすう}{クーロンの定数}\bkidx{しんくうのゆうでんりつ}{真空の誘電率}F/m（ファラド毎メートル）は、電-3 で学ぶコンデンサーの単位 F を使った書き方で、N·m²/C² の逆数と同じ単位です。

\begin{phyomi}[$F=k\,\dfrac{q_1q_2}{r^2}$]
「電気の力は、2つの電荷のかけ算に比例し、距離の2乗に反比例する」。$q_1q_2$ が正（同じ符号）なら、はね返す力。負（違う符号）なら、引き合う力。重力の $F=GMm/r^2$（力-2）とそっくりの形です。違いは、(1) 電荷には正と負があって、はね返す力もあること、(2) $k$ が $G$ の $10^{20}$ 倍も大きいことです。
\end{phyomi}

\noindent なぜ $k$ を、わざわざ $4\pi$ と $\varepsilon_0$ で書くのでしょう。$4\pi r^2$ は半径 $r$ の球の表面積です。電荷から出た「電気の影響」が、球面に広がって薄まっていくと考えると、$4\pi$ がつくのは自然です。このことは、電-2 のガウスの法則ではっきりします。電-3 では、物の中の電気の力を考えるのに、$\varepsilon_0$ のほうが便利になります。

\begin{phbaisuru}{距離 $r$}
力は $\tfrac14$ になります。3倍にすると $\tfrac19$、10倍にすると $\tfrac{1}{100}$ です。下のグラフは、\qty{1}{\micro C}どうしの力です（\unit{\micro C}（マイクロクーロン）は $10^{-6}$\,C）。距離が近いと、力は急に大きくなります。
\par\smallskip\noindent\begin{tikzpicture}
\begin{axis}[bkplot, width=100mm, height=44mm, xmin=0, xmax=0.33, ymin=0, ymax=4.6, xtick={0.05,0.1,0.2,0.3}, xticklabel style={/pgf/number format/fixed}, ytick={0.9,3.6}, yticklabels={0.90,3.6}, xlabel={$r$（m）}, ylabel={$F$（N）}]
  \addplot[bkcurve, domain=0.045:0.32, samples=80] {8.99e9*1e-12/x^2};
  \fill[bkRed] (axis cs:0.05,3.596) circle (1.8pt) (axis cs:0.1,0.899) circle (1.8pt) (axis cs:0.15,0.3996) circle (1.8pt);
  \node[font=\sffamily\footnotesize, anchor=west, text=bkRed!75!black] at (axis cs:0.11,1.6) {距離2倍で $\tfrac14$、3倍で $\tfrac19$};
\end{axis}
\end{tikzpicture}
\end{phbaisuru}

\begin{phbaisuru}{電荷 $q_1$}
力も2倍になります。$q_1$ と $q_2$ の両方を2倍にすると4倍です。式が $q_1q_2$ のかけ算だからです。
\end{phbaisuru}

\begin{phmijika}[1\,C は、とても大きな電荷]
\qty{1}{C}どうしを\qty{1}{m}離して置くと、力は約 $9\times10^9$\,N で、約90万トンの物の重さと同じです（例題1）。ふつうの静電気の電荷は、\qty{1}{nC}から\qty{1}{\micro C}くらいで、1\,C よりずっと小さいのです。雷の1回の放電で、雲から地面へ運ばれる電荷でも、数 C から数十 C ほどです。いっぽう、電流として流れる電荷は、もっとずっと多く、スマホの電池（\qty{3000}{mAh}）が1回の充電で送り出す電荷は約 \num{10800}\,C です（コラム）。電流では、正と負の電荷がほとんど同じだけあって力が打ち消されるので、こんなに大きな電荷が動いても、ばらばらになりません。
\end{phmijika}

\begin{seikaku}
クーロンの法則は、実験から見つけた法則で、ほかの式から導いたものではありません。電-2 で学ぶガウスの法則と、内容は同じです。また、$k$（と $\varepsilon_0$）の値は、2019年の SI の改定から、測って決める値になりました。$\varepsilon_0=\qty{8.8541878188e-12}{F/m}$（CODATA 2022）で、$k=\qty{8.98755e9}{N.m^2/C^2}$ です。この本では、数学の本（8-5⑥）と同じく、$k=8.99\times10^9$、$\varepsilon_0=8.85\times10^{-12}$ を使います。
\end{seikaku}

\phboxhead{矢印で書く}
力は向きをもつ量（ベクトル）です。電荷 $q_1$ から $q_2$ へ向かう長さ1の矢印を $\vec e_{12}$ とすると、$q_2$ が受ける力は
\pheq{\vec F_{2}=k\,\frac{q_1q_2}{r^2}\,\vec e_{12}}{N}
と書けます。$q_1q_2>0$ なら $\vec e_{12}$ の向き（$q_1$ から遠ざかる向き）、$q_1q_2<0$ なら反対の向き（$q_1$ へ近づく向き）です。$q_1$ が受ける力は、ちょうど反対向きの $-\vec F_2$ です。

\begin{phmodel}[点電荷]
\phignore{電荷をもつ物の大きさと形。周りの空気。}
\phassume{電荷は1つの点に集まっている（\textbf{点電荷}）。電荷は止まっている。}
\phlimit{物の大きさが、距離と同じくらいになるとき（⑦のように、広がりを積分する）。物が水などの中にあるとき（水の中では力が約80分の1：電-3）。電荷が速く動くとき（磁気の力が加わる：電-6）。}
\end{phmodel}

\phboxhead{はてな？の答え}
\label{ph:電-1:hatena}\bkhatenaref の答えです。頭をこすった下じきは、髪の毛から電子をもらって負に帯電し、髪の毛は正に帯電しています。下じきを持ち上げると、髪の毛の先の正の電荷が、下じきの負の電荷に引かれます（違う符号は引き合う）。例題3で見積もると、髪の毛の先に $10^{-10}$\,C くらいの、ほんのわずかな電荷があるだけで、電気の力は髪の毛1本の重さと同じくらいになります。だから、髪の毛は重力に逆らって立ち上がるのです。さらに、髪の毛どうしはどれも正に帯電しているので、たがいにはね返し合い、ばらばらに広がります。下じきを遠ざけると、力は距離の2乗に反比例して弱くなり、髪は重力に負けて下がります（じぶんでとこう3）。

%% ---------------------------------------------------------------
\bksubsection{3}{重ね合わせ：力は矢印でたす}
\label{ph:電-1:sec3}
電荷が3つ以上あるときは、どうなるでしょう。たとえば、電荷 $q$ の周りに $Q_1$ と $Q_2$ があるとき、$q$ が受ける力は、$Q_1$ だけがあるときの力 $\vec F_1$ と、$Q_2$ だけがあるときの力 $\vec F_2$ を、\textbf{矢印としてたしたもの}（ベクトルの和：4-10③）になります。
\pheq{\vec F=\vec F_1+\vec F_2+\cdots}{N}
これを\textbf{重ね合わせの原理}といいます。\bkidx{かさねあわせのげんり}{重ね合わせの原理}$Q_2$ があっても、$Q_1$ と $q$ のあいだの力は変わらない、ということです。実験で確かめられている、電気の力の大切な性質です。

\begin{bkfigure}
\begin{tikzpicture}[x=1mm, y=1mm]
  \node[sxb lab, anchor=north] at (0,-3.4) {$Q_1$（＋）};
  \node[sxb lab, anchor=north] at (50,-3.4) {$Q_2$（－）};
  \node[sxb lab, anchor=south east] at (23.4,26.4) {$q$（＋）};
  \draw[bkGray, line width=0.4pt, dashed] (0,0) -- (25,25) -- (50,0);
  \draw[bkRed!80!black, sxb f] (25,25) -- ++(45:14) coordinate (F1);
  \draw[bkBlue!80!black, sxb f] (25,25) -- ++(-45:14) coordinate (F2);
  \draw[bkGray, line width=0.5pt, dashed] (F1) -- ++(-45:14) coordinate (S);
  \draw[bkGray, line width=0.5pt, dashed] (F2) -- (S);
  \draw[bkInk, line width=1.6pt, -{Stealth[length=3mm, width=2.4mm]}] (25,25) -- (S);
  \sxbplus{(25,25)}{2.2}\sxbplus{(0,0)}{2.4}\sxbminus{(50,0)}{2.4}
  \node[sxb lab, anchor=south east, text=bkRed!75!black] at (F1) {$\vec F_1$};
  \node[sxb lab, anchor=north east, text=bkBlue!75!black] at (F2) {$\vec F_2$};
  \node[sxb lab, anchor=west] at ($(S)+(1,0)$) {$\vec F=\vec F_1+\vec F_2$};
  \node[anchor=west, font=\small, text width=44mm, align=left] at (70,12) {%
  $Q_1$ は $q$ をはね返し（$Q_1$ から遠ざかる向き）、$Q_2$ は $q$ を引きつける（$Q_2$ へ向かう向き）。\par\vspace{2pt}
  2つの矢印を、平行四辺形の対角線のようにたす。};
\end{tikzpicture}
\bkfigcaption{重ね合わせ。$q$ が受ける力は、$Q_1$ からの力と $Q_2$ からの力の、ベクトルとしての和。この図では、たての成分が打ち消されて、右向きになる（例題4）。}\label{ph:電-1:fig-kasane}
\end{bkfigure}

\noindent 計算では、それぞれの力を $x$ 成分と $y$ 成分に分け、成分ごとにたします（4-10）。対称な置き方では、成分の一部が打ち消し合って、計算が楽になります（例題4・6）。

%% ---------------------------------------------------------------
\Needspace{12\baselineskip}\bksubsection{4}{電場：電荷の周りの空間の性質}
\label{ph:電-1:sec4}
重ね合わせの計算では、力を受ける電荷 $q$ を置く場所を変えるたびに、全部の力を計算し直さなければなりません。そこで、考え方を少し変えます。

\noindent\begin{minipage}[c]{0.53\linewidth}\raggedright
電荷 $Q$ があると、その周りの空間の\textbf{どの場所にも}、「そこに電荷を置いたら、1\,C あたりこれだけの力がはたらく」という矢印が決まっている、と考えるのです。この矢印を\textbf{電場}（電界）といい、$\vec E$ と書きます。\bkidx{でんば}{電場}

ある場所に小さな電荷 $q$（\textbf{試しの電荷}）を置いて、受ける力が $\vec F$ なら、
\end{minipage}\hfill
\begin{minipage}[c]{0.44\linewidth}\centering
\begin{tikzpicture}[x=0.85mm, y=0.85mm]
  \sxbplus{(0,0)}{2.6}
  \node[sxb lab, inner sep=0pt] at (112.5:8) {$Q$};
  \foreach \a in {0,45,...,315} {\draw[sxb arr, line width=0.9pt] (\a:5) -- (\a:13);
    \draw[sxb arr] (\a:20) -- (\a:22);}
  \filldraw[fill=bkGreen!40, draw=bkInk, line width=0.4pt] (30,0) circle (1.2);
  \draw[sxb arr] (31.5,0) -- (35,0);
  \node[sxb lab, anchor=north] at (30,-2.4) {試しの電荷 $q$};
  \node[sxb lab, anchor=south, align=center] at (16,16) {近くは強く、\\遠くは弱い};
\end{tikzpicture}
\end{minipage}
\pheq{\vec E=\frac{\vec F}{q}\qquad\text{（$\vec F=q\vec E$）}}{N/C ＝ N÷C}
\label{ph:電-1:tenkaba}です。電場の単位は N/C です。電-2 で学ぶ電位の単位 V（ボルト）を使うと V/m と書け、N/C と V/m は同じ単位です。正の電荷は電場の向きに、負の電荷は電場と反対の向きに力を受けます。

\phboxhead{点電荷の電場}
電荷 $Q$ から距離 $r$ の所に $q$ を置くと、クーロンの法則で、力の大きさは $kQq/r^2$ です。これを $q$ でわると、
\pheq{E=k\,\frac{Q}{r^2}}{N/C ＝ $\dfrac{\text{N·m}^2}{\text{C}^2}$×$\dfrac{\text{C}}{\text{m}^2}$}
です。向きは、$Q$ が正なら $Q$ から外向き、負なら $Q$ へ向かう向きです。試しの電荷 $q$ が消えて、電場が $Q$ と場所 $r$ だけで決まっていることに注意しましょう。これが「空間の性質」と考えられる理由です。\bkidx{てんでんか}{点電荷}

\begin{phflow}[クーロンの法則から点電荷の電場へ]
\phflowrow{$\vec F=k\,\dfrac{Qq}{r^2}\,\vec e_r$}{試しの電荷 $q$ で\\わる}{$\vec E=k\,\dfrac{Q}{r^2}\,\vec e_r$}
\phflowrow{$\vec E=\sum k\,\dfrac{Q_i}{r_i^{\,2}}\,\vec e_i$}{電荷が1つだけ\\大きさだけ見る}{$E=k\,\dfrac{Q}{r^2}$}
\end{phflow}
\noindent $\vec e_r$ は、$Q$ からその場所へ向かう長さ1の矢印です。電荷がいくつもあるときの電場は、③の重ね合わせで、それぞれの電荷がつくる電場の矢印の和です（2段めの左の式の $\sum$ は「全部たす」の記号：5-10）。

\begin{bkcenter}
\begin{bktable}{colspec={Q[c,m,wd=30mm] *{5}{Q[c,m,wd=15mm]}}, rowsep=1.5pt, row{1}={bg=bkTheme!10, font=\sffamily\small}}
距離 $r$〔m〕 & 0.1 & 0.2 & 0.3 & 0.5 & 1 \\
電場 $E$〔N/C〕 & 899 & 225 & 99.9 & 36.0 & 8.99 \\
\end{bktable}
\bktabcaption{\qty{1}{nC}の点電荷の電場（8-5⑥ の表と同じ）。距離が2倍で $\tfrac14$、10倍で $\tfrac1{100}$。}\label{ph:電-1:tab-1nC}
\end{bkcenter}

\begin{phyomi}[$E=k\,\dfrac{Q}{r^2}$]
「電荷が大きいほど、近いほど、電場は強い。距離の2乗で弱まる」。表\ref{ph:電-1:tab-1nC}の\qty{0.1}{m}の所に、\qty{1}{nC}の試しの電荷を置くと、受ける力は $899\times10^{-9}\approx\qty{9.0e-7}{N}$ です。電子（電荷 $-e$）を置くと、電場と反対向きに $899\times1.60\times10^{-19}\approx\qty{1.4e-16}{N}$ の力を受けます。
\end{phyomi}

\begin{phmijika}[なぜ「場」と考えるのか]
2つの電荷のあいだの力を、「離れた物どうしが、何もない空間を飛びこえて、いきなり引き合う」と考えるのは、少し不思議です。19世紀のイギリスのファラデーは、電荷の周りの空間そのものが、何かの状態（場）をもっていて、それが別の電荷に力をおよぼす、と考えました。この見方は、ただの言いかえではありません。電荷をゆらすと、そのゆれは、周りの電場をゆらしながら光の速さで伝わっていきます。それが電波や光です（電-15）。電場はエネルギーももっています（電-3）。
\end{phmijika}

\begin{phbaisuru}{電荷 $Q$}
どの場所でも電場が2倍になります。電場の地図の矢印が、全部2倍に伸びると考えればよいのです。いっぽう、距離を2倍にすると $\tfrac14$ になります（グラフ。\qty{1}{nC}と\qty{2}{nC}）。
\par\smallskip\noindent\begin{tikzpicture}
\begin{axis}[bkplot, width=100mm, height=44mm, xmin=0, xmax=0.55, ymin=0, ymax=2100, xtick={0.1,0.2,0.3,0.4,0.5}, xticklabel style={/pgf/number format/fixed}, ytick={225,899,1798}, xlabel={$r$（m）}, ylabel={$E$（N/C）}]
  \addplot[bkcurve, domain=0.066:0.52, samples=80] {8.99/x^2};
  \addplot[bkOrange, line width=1.6pt, dash pattern=on 4pt off 2pt, domain=0.093:0.52, samples=80] {2*8.99/x^2};
  \fill[bkRed] (axis cs:0.1,899) circle (1.8pt) (axis cs:0.2,224.75) circle (1.8pt) (axis cs:0.1,1798) circle (1.8pt);
  \node[font=\sffamily\footnotesize, anchor=west] at (axis cs:0.2,700) {\qty{2}{nC}（破線）};
  \node[font=\sffamily\footnotesize, anchor=west] at (axis cs:0.3,330) {\qty{1}{nC}（実線）};
\end{axis}
\end{tikzpicture}
\end{phbaisuru}


%% ---------------------------------------------------------------
\Needspace{12\baselineskip}\bksubsection{5}{電気力線：電場を線で見る}
\label{ph:電-1:sec5}
電場の矢印を、場所ごとにかくと、矢印だらけになって見にくくなります。そこで、8-5④ の力線のように、矢印の向きをつないだ線をかきます。これを\textbf{電気力線}といいます。\bkidx{でんきりきせん}{電気力線}電気力線には、次の決まりがあります。
\begin{enumerate}[label={\bkcircled{\arabic*}}, leftmargin=2.4\zw, itemsep=1pt, topsep=2pt]
\item 正の電荷から出て、負の電荷に入る（または、遠くへ出ていく・遠くから来る）。
\item 線のどの点でも、線の向き（接線の向き）が、その点の電場の向き。
\item 線がこみ合っている所ほど、電場が強い。
\item 電気力線は、たがいに交わらない（1つの場所の電場の向きは1つだから）。
\item 電荷が大きいほど、たくさんの線が出る（または入る）。
\end{enumerate}

\begin{bkfigure}
\begin{tikzpicture}[x=0.6mm, y=0.6mm]
\begin{scope}
  \clip (-52,-30) rectangle (52,30);
\draw[sxb line] plot[smooth] coordinates {(-15.0,0.0) (-12.8,0.0) (-11.3,0.0) (-9.8,0.0) (-8.3,0.0) (-6.8,0.0) (-5.3,0.0) (-3.8,0.0) (-2.3,0.0) (-0.8,0.0) (0.7,0.0) (2.2,0.0) (3.7,0.0) (5.2,0.0) (6.7,0.0) (8.2,0.0) (9.7,0.0) (11.2,0.0) (12.7,0.0) (15.0,0.0)}; \draw[sxb tip] (0.40,0.00) -- (0.90,0.00);
\draw[sxb line] plot[smooth] coordinates {(-15.0,0.0) (-13.1,1.1) (-11.8,1.8) (-10.5,2.6) (-9.1,3.3) (-7.8,3.9) (-6.4,4.5) (-5.0,5.0) (-3.6,5.4) (-2.1,5.7) (-0.6,5.8) (0.9,5.8) (2.4,5.6) (3.9,5.3) (5.3,4.9) (6.7,4.4) (8.1,3.8) (9.4,3.1) (10.8,2.4) (12.1,1.7) (13.2,1.1) (15.0,0.0)}; \draw[sxb tip] (-0.90,5.80) -- (-0.40,5.80);
\draw[sxb line] plot[smooth] coordinates {(-15.0,0.0) (-13.9,1.9) (-13.1,3.2) (-12.3,4.5) (-11.5,5.7) (-10.6,6.9) (-9.7,8.1) (-8.7,9.2) (-7.6,10.2) (-6.4,11.1) (-5.1,11.9) (-3.8,12.6) (-2.3,13.0) (-0.8,13.2) (0.7,13.3) (2.1,13.0) (3.6,12.6) (4.9,12.0) (6.2,11.3) (7.4,10.4) (8.5,9.3) (9.6,8.3) (10.5,7.1) (11.4,5.9) (12.2,4.6) (13.0,3.4) (13.8,2.1) (15.0,0.0)}; \draw[sxb tip] (0.41,13.36) -- (0.90,13.26);
\draw[sxb line] plot[smooth] coordinates {(-15.0,0.0) (-15.0,2.2) (-15.0,3.7) (-14.9,5.2) (-14.9,6.7) (-14.7,8.2) (-14.5,9.7) (-14.3,11.2) (-14.0,12.6) (-13.6,14.1) (-13.1,15.5) (-12.5,16.9) (-11.9,18.2) (-11.1,19.5) (-10.2,20.7) (-9.2,21.9) (-8.2,22.9) (-7.0,23.8) (-5.7,24.6) (-4.3,25.2) (-2.9,25.7) (-1.4,26.0) (0.1,26.1) (1.6,26.0) (3.0,25.7) (4.4,25.2) (5.8,24.6) (7.1,23.8) (8.2,22.8) (9.3,21.8) (10.3,20.6) (11.2,19.4) (11.9,18.1) (12.6,16.8) (13.1,15.4) (13.6,14.0) (14.0,12.5) (14.3,11.0) (14.6,9.6) (14.7,8.1) (14.9,6.6) (14.9,5.1) (15.0,3.6) (15.0,2.1) (15.0,0.0)}; \draw[sxb tip] (-0.20,26.12) -- (0.30,26.09);
\draw[sxb line] plot[smooth] coordinates {(-15.0,0.0) (-16.1,1.9) (-16.8,3.2) (-17.6,4.5) (-18.3,5.9) (-18.9,7.2) (-19.6,8.5) (-20.2,9.9) (-20.8,11.3) (-21.3,12.7) (-21.8,14.1) (-22.3,15.5) (-22.7,17.0) (-23.1,18.4) (-23.4,19.9) (-23.7,21.4) (-23.9,22.8) (-24.1,24.3) (-24.3,25.8) (-24.3,27.3) (-24.3,29.8)}; \draw[sxb tip] (-21.20,12.42) -- (-21.37,12.89);
\draw[sxb line] plot[smooth] coordinates {(-15.0,0.0) (-16.9,1.1) (-18.2,1.9) (-19.5,2.6) (-20.8,3.4) (-22.1,4.2) (-23.3,5.0) (-24.6,5.8) (-25.8,6.6) (-27.1,7.4) (-28.3,8.3) (-29.5,9.2) (-30.8,10.0) (-32.0,10.9) (-33.2,11.8) (-34.3,12.8) (-35.5,13.7) (-36.7,14.6) (-37.8,15.6) (-39.0,16.6) (-40.1,17.6) (-41.2,18.6) (-42.3,19.6) (-43.4,20.6) (-44.5,21.6) (-45.6,22.7) (-46.6,23.8) (-47.7,24.8) (-48.7,25.9) (-49.7,27.0) (-50.8,28.1) (-51.9,29.4)}; \draw[sxb tip] (-32.98,11.60) -- (-33.35,11.93);
\draw[sxb line] plot[smooth] coordinates {(-15.0,0.0) (-17.2,0.0) (-18.7,0.0) (-20.2,0.0) (-21.7,0.0) (-23.2,0.0) (-24.7,0.0) (-26.2,0.0) (-27.7,0.0) (-29.2,0.0) (-30.7,0.0) (-32.2,0.0) (-33.7,0.0) (-35.2,0.0) (-36.7,0.0) (-38.2,0.0) (-39.7,0.0) (-41.2,0.0) (-42.7,0.0) (-44.2,0.0) (-45.7,0.0) (-47.2,0.0) (-48.7,0.0) (-50.2,0.0) (-52.0,0.0)}; \draw[sxb tip] (-31.90,0.00) -- (-32.40,0.00);
\draw[sxb line] plot[smooth] coordinates {(-15.0,0.0) (-16.9,-1.1) (-18.2,-1.9) (-19.5,-2.6) (-20.8,-3.4) (-22.1,-4.2) (-23.3,-5.0) (-24.6,-5.8) (-25.8,-6.6) (-27.1,-7.4) (-28.3,-8.3) (-29.5,-9.2) (-30.8,-10.0) (-32.0,-10.9) (-33.2,-11.8) (-34.3,-12.8) (-35.5,-13.7) (-36.7,-14.6) (-37.8,-15.6) (-39.0,-16.6) (-40.1,-17.6) (-41.2,-18.6) (-42.3,-19.6) (-43.4,-20.6) (-44.5,-21.6) (-45.6,-22.7) (-46.6,-23.8) (-47.7,-24.8) (-48.7,-25.9) (-49.7,-27.0) (-50.8,-28.1) (-51.9,-29.4)}; \draw[sxb tip] (-32.98,-11.60) -- (-33.35,-11.93);
\draw[sxb line] plot[smooth] coordinates {(-15.0,0.0) (-16.1,-1.9) (-16.8,-3.2) (-17.6,-4.5) (-18.3,-5.9) (-18.9,-7.2) (-19.6,-8.5) (-20.2,-9.9) (-20.8,-11.3) (-21.3,-12.7) (-21.8,-14.1) (-22.3,-15.5) (-22.7,-17.0) (-23.1,-18.4) (-23.4,-19.9) (-23.7,-21.4) (-23.9,-22.8) (-24.1,-24.3) (-24.3,-25.8) (-24.3,-27.3) (-24.3,-29.8)}; \draw[sxb tip] (-21.20,-12.42) -- (-21.37,-12.89);
\draw[sxb line] plot[smooth] coordinates {(-15.0,0.0) (-15.0,-2.2) (-15.0,-3.7) (-14.9,-5.2) (-14.9,-6.7) (-14.7,-8.2) (-14.5,-9.7) (-14.3,-11.2) (-14.0,-12.6) (-13.6,-14.1) (-13.1,-15.5) (-12.5,-16.9) (-11.9,-18.2) (-11.1,-19.5) (-10.2,-20.7) (-9.2,-21.9) (-8.2,-22.9) (-7.0,-23.8) (-5.7,-24.6) (-4.3,-25.2) (-2.9,-25.7) (-1.4,-26.0) (0.1,-26.1) (1.6,-26.0) (3.0,-25.7) (4.4,-25.2) (5.8,-24.6) (7.1,-23.8) (8.2,-22.8) (9.3,-21.8) (10.3,-20.6) (11.2,-19.4) (11.9,-18.1) (12.6,-16.8) (13.1,-15.4) (13.6,-14.0) (14.0,-12.5) (14.3,-11.0) (14.6,-9.6) (14.7,-8.1) (14.9,-6.6) (14.9,-5.1) (15.0,-3.6) (15.0,-2.1) (15.0,0.0)}; \draw[sxb tip] (-0.20,-26.12) -- (0.30,-26.09);
\draw[sxb line] plot[smooth] coordinates {(-15.0,0.0) (-13.9,-1.9) (-13.1,-3.2) (-12.3,-4.5) (-11.5,-5.7) (-10.6,-6.9) (-9.7,-8.1) (-8.7,-9.2) (-7.6,-10.2) (-6.4,-11.1) (-5.1,-11.9) (-3.8,-12.6) (-2.3,-13.0) (-0.8,-13.2) (0.7,-13.3) (2.1,-13.0) (3.6,-12.6) (4.9,-12.0) (6.2,-11.3) (7.4,-10.4) (8.5,-9.3) (9.6,-8.3) (10.5,-7.1) (11.4,-5.9) (12.2,-4.6) (13.0,-3.4) (13.8,-2.1) (15.0,0.0)}; \draw[sxb tip] (0.41,-13.36) -- (0.90,-13.26);
\draw[sxb line] plot[smooth] coordinates {(-15.0,0.0) (-13.1,-1.1) (-11.8,-1.8) (-10.5,-2.6) (-9.1,-3.3) (-7.8,-3.9) (-6.4,-4.5) (-5.0,-5.0) (-3.6,-5.4) (-2.1,-5.7) (-0.6,-5.8) (0.9,-5.8) (2.4,-5.6) (3.9,-5.3) (5.3,-4.9) (6.7,-4.4) (8.1,-3.8) (9.4,-3.1) (10.8,-2.4) (12.1,-1.7) (13.2,-1.1) (15.0,0.0)}; \draw[sxb tip] (-0.90,-5.80) -- (-0.40,-5.80);
\end{scope}
\draw[bkGray, line width=0.4pt] (-52,-30) rectangle (52,30);
\sxbplus{(-15,0)}{3.2}\sxbminus{(15,0)}{3.2}
\node[sxb lab, anchor=north] at (0,-31) {正と負（同じ大きさ）};
\begin{scope}[shift={(112,0)}]
  \begin{scope}
  \clip (-52,-30) rectangle (52,30);
\draw[sxb line] plot[smooth] coordinates {(-15.0,0.0) (-13.1,1.1) (-11.8,1.9) (-10.5,2.6) (-9.3,3.5) (-8.0,4.3) (-6.9,5.3) (-5.8,6.3) (-4.9,7.5) (-4.2,8.8) (-3.6,10.2) (-3.3,11.7) (-3.0,13.2) (-2.8,14.6) (-2.7,16.1) (-2.6,17.6) (-2.6,19.1) (-2.6,20.6) (-2.6,22.1) (-2.6,23.6) (-2.6,25.1) (-2.7,26.6) (-2.7,28.1) (-2.8,29.9)}; \draw[sxb tip] (-3.66,9.91) -- (-3.56,10.40);
\draw[sxb line] plot[smooth] coordinates {(-15.0,0.0) (-13.9,1.9) (-13.2,3.2) (-12.5,4.5) (-11.8,5.9) (-11.2,7.3) (-10.6,8.7) (-10.2,10.1) (-9.8,11.5) (-9.6,13.0) (-9.4,14.5) (-9.3,16.0) (-9.2,17.5) (-9.3,19.0) (-9.3,20.5) (-9.4,22.0) (-9.6,23.5) (-9.7,25.0) (-9.9,26.5) (-10.1,27.9) (-10.5,29.9)}; \draw[sxb tip] (-9.64,12.70) -- (-9.57,13.20);
\draw[sxb line] plot[smooth] coordinates {(-15.0,0.0) (-15.0,2.2) (-15.0,3.7) (-15.1,5.2) (-15.1,6.7) (-15.3,8.2) (-15.4,9.7) (-15.6,11.2) (-15.9,12.6) (-16.2,14.1) (-16.6,15.6) (-17.0,17.0) (-17.4,18.5) (-17.8,19.9) (-18.3,21.3) (-18.8,22.7) (-19.4,24.1) (-19.9,25.5) (-20.5,26.9) (-21.1,28.3) (-21.8,29.9)}; \draw[sxb tip] (-16.12,13.81) -- (-16.25,14.29);
\draw[sxb line] plot[smooth] coordinates {(-15.0,0.0) (-16.1,1.9) (-16.9,3.2) (-17.6,4.5) (-18.4,5.8) (-19.3,7.0) (-20.1,8.3) (-20.9,9.5) (-21.8,10.7) (-22.7,11.9) (-23.6,13.1) (-24.6,14.3) (-25.5,15.4) (-26.5,16.6) (-27.5,17.7) (-28.4,18.8) (-29.4,20.0) (-30.4,21.1) (-31.5,22.2) (-32.5,23.3) (-33.5,24.4) (-34.5,25.5) (-35.6,26.5) (-36.6,27.6) (-37.7,28.7) (-38.9,29.9)}; \draw[sxb tip] (-24.41,14.07) -- (-24.73,14.45);
\draw[sxb line] plot[smooth] coordinates {(-15.0,0.0) (-16.9,1.1) (-18.2,1.8) (-19.5,2.6) (-20.8,3.3) (-22.1,4.0) (-23.5,4.7) (-24.8,5.4) (-26.1,6.1) (-27.5,6.8) (-28.8,7.5) (-30.2,8.1) (-31.5,8.8) (-32.9,9.4) (-34.2,10.1) (-35.6,10.7) (-36.9,11.3) (-38.3,12.0) (-39.7,12.6) (-41.0,13.2) (-42.4,13.8) (-43.8,14.4) (-45.1,15.0) (-46.5,15.6) (-47.9,16.2) (-49.3,16.8) (-50.7,17.4) (-51.8,17.9)}; \draw[sxb tip] (-31.22,8.68) -- (-31.68,8.88);
\draw[sxb line] plot[smooth] coordinates {(-15.0,0.0) (-17.2,0.0) (-18.7,0.0) (-20.2,0.0) (-21.7,0.0) (-23.2,0.0) (-24.7,0.0) (-26.2,0.0) (-27.7,0.0) (-29.2,0.0) (-30.7,0.0) (-32.2,0.0) (-33.7,0.0) (-35.2,0.0) (-36.7,0.0) (-38.2,0.0) (-39.7,0.0) (-41.2,0.0) (-42.7,0.0) (-44.2,0.0) (-45.7,0.0) (-47.2,0.0) (-48.7,0.0) (-50.2,0.0) (-52.0,0.0)}; \draw[sxb tip] (-31.90,0.00) -- (-32.40,0.00);
\draw[sxb line] plot[smooth] coordinates {(-15.0,0.0) (-16.9,-1.1) (-18.2,-1.8) (-19.5,-2.6) (-20.8,-3.3) (-22.1,-4.0) (-23.5,-4.7) (-24.8,-5.4) (-26.1,-6.1) (-27.5,-6.8) (-28.8,-7.5) (-30.2,-8.1) (-31.5,-8.8) (-32.9,-9.4) (-34.2,-10.1) (-35.6,-10.7) (-36.9,-11.3) (-38.3,-12.0) (-39.7,-12.6) (-41.0,-13.2) (-42.4,-13.8) (-43.8,-14.4) (-45.1,-15.0) (-46.5,-15.6) (-47.9,-16.2) (-49.3,-16.8) (-50.7,-17.4) (-51.8,-17.9)}; \draw[sxb tip] (-31.22,-8.68) -- (-31.68,-8.88);
\draw[sxb line] plot[smooth] coordinates {(-15.0,0.0) (-16.1,-1.9) (-16.9,-3.2) (-17.6,-4.5) (-18.4,-5.8) (-19.3,-7.0) (-20.1,-8.3) (-20.9,-9.5) (-21.8,-10.7) (-22.7,-11.9) (-23.6,-13.1) (-24.6,-14.3) (-25.5,-15.4) (-26.5,-16.6) (-27.5,-17.7) (-28.4,-18.8) (-29.4,-20.0) (-30.4,-21.1) (-31.5,-22.2) (-32.5,-23.3) (-33.5,-24.4) (-34.5,-25.5) (-35.6,-26.5) (-36.6,-27.6) (-37.7,-28.7) (-38.9,-29.9)}; \draw[sxb tip] (-24.41,-14.07) -- (-24.73,-14.45);
\draw[sxb line] plot[smooth] coordinates {(-15.0,0.0) (-15.0,-2.2) (-15.0,-3.7) (-15.1,-5.2) (-15.1,-6.7) (-15.3,-8.2) (-15.4,-9.7) (-15.6,-11.2) (-15.9,-12.6) (-16.2,-14.1) (-16.6,-15.6) (-17.0,-17.0) (-17.4,-18.5) (-17.8,-19.9) (-18.3,-21.3) (-18.8,-22.7) (-19.4,-24.1) (-19.9,-25.5) (-20.5,-26.9) (-21.1,-28.3) (-21.8,-29.9)}; \draw[sxb tip] (-16.12,-13.81) -- (-16.25,-14.29);
\draw[sxb line] plot[smooth] coordinates {(-15.0,0.0) (-13.9,-1.9) (-13.2,-3.2) (-12.5,-4.5) (-11.8,-5.9) (-11.2,-7.3) (-10.6,-8.7) (-10.2,-10.1) (-9.8,-11.5) (-9.6,-13.0) (-9.4,-14.5) (-9.3,-16.0) (-9.2,-17.5) (-9.3,-19.0) (-9.3,-20.5) (-9.4,-22.0) (-9.6,-23.5) (-9.7,-25.0) (-9.9,-26.5) (-10.1,-27.9) (-10.5,-29.9)}; \draw[sxb tip] (-9.64,-12.70) -- (-9.57,-13.20);
\draw[sxb line] plot[smooth] coordinates {(-15.0,0.0) (-13.1,-1.1) (-11.8,-1.9) (-10.5,-2.6) (-9.3,-3.5) (-8.0,-4.3) (-6.9,-5.3) (-5.8,-6.3) (-4.9,-7.5) (-4.2,-8.8) (-3.6,-10.2) (-3.3,-11.7) (-3.0,-13.2) (-2.8,-14.6) (-2.7,-16.1) (-2.6,-17.6) (-2.6,-19.1) (-2.6,-20.6) (-2.6,-22.1) (-2.6,-23.6) (-2.6,-25.1) (-2.7,-26.6) (-2.7,-28.1) (-2.8,-29.9)}; \draw[sxb tip] (-3.66,-9.91) -- (-3.56,-10.40);
\draw[sxb line] plot[smooth] coordinates {(15.0,0.0) (17.2,0.0) (18.7,0.0) (20.2,0.0) (21.7,0.0) (23.2,0.0) (24.7,0.0) (26.2,0.0) (27.7,0.0) (29.2,0.0) (30.7,0.0) (32.2,0.0) (33.7,0.0) (35.2,0.0) (36.7,0.0) (38.2,0.0) (39.7,0.0) (41.2,0.0) (42.7,0.0) (44.2,0.0) (45.7,0.0) (47.2,0.0) (48.7,0.0) (50.2,0.0) (52.0,0.0)}; \draw[sxb tip] (31.90,0.00) -- (32.40,0.00);
\draw[sxb line] plot[smooth] coordinates {(15.0,0.0) (16.9,1.1) (18.2,1.8) (19.5,2.6) (20.8,3.3) (22.1,4.0) (23.5,4.7) (24.8,5.4) (26.1,6.1) (27.5,6.8) (28.8,7.5) (30.2,8.1) (31.5,8.8) (32.9,9.4) (34.2,10.1) (35.6,10.7) (36.9,11.3) (38.3,12.0) (39.7,12.6) (41.0,13.2) (42.4,13.8) (43.8,14.4) (45.1,15.0) (46.5,15.6) (47.9,16.2) (49.3,16.8) (50.7,17.4) (51.8,17.9)}; \draw[sxb tip] (31.22,8.68) -- (31.68,8.88);
\draw[sxb line] plot[smooth] coordinates {(15.0,0.0) (16.1,1.9) (16.9,3.2) (17.6,4.5) (18.4,5.8) (19.3,7.0) (20.1,8.3) (20.9,9.5) (21.8,10.7) (22.7,11.9) (23.6,13.1) (24.6,14.3) (25.5,15.4) (26.5,16.6) (27.5,17.7) (28.4,18.8) (29.4,20.0) (30.4,21.1) (31.5,22.2) (32.5,23.3) (33.5,24.4) (34.5,25.5) (35.6,26.5) (36.6,27.6) (37.7,28.7) (38.9,29.9)}; \draw[sxb tip] (24.41,14.07) -- (24.73,14.45);
\draw[sxb line] plot[smooth] coordinates {(15.0,0.0) (15.0,2.2) (15.0,3.7) (15.1,5.2) (15.1,6.7) (15.3,8.2) (15.4,9.7) (15.6,11.2) (15.9,12.6) (16.2,14.1) (16.6,15.6) (17.0,17.0) (17.4,18.5) (17.8,19.9) (18.3,21.3) (18.8,22.7) (19.4,24.1) (19.9,25.5) (20.5,26.9) (21.1,28.3) (21.8,29.9)}; \draw[sxb tip] (16.12,13.81) -- (16.25,14.29);
\draw[sxb line] plot[smooth] coordinates {(15.0,0.0) (13.9,1.9) (13.2,3.2) (12.5,4.5) (11.8,5.9) (11.2,7.3) (10.6,8.7) (10.2,10.1) (9.8,11.5) (9.6,13.0) (9.4,14.5) (9.3,16.0) (9.2,17.5) (9.3,19.0) (9.3,20.5) (9.4,22.0) (9.6,23.5) (9.7,25.0) (9.9,26.5) (10.1,27.9) (10.5,29.9)}; \draw[sxb tip] (9.64,12.70) -- (9.57,13.20);
\draw[sxb line] plot[smooth] coordinates {(15.0,0.0) (13.1,1.1) (11.8,1.9) (10.5,2.6) (9.3,3.5) (8.0,4.3) (6.9,5.3) (5.8,6.3) (4.9,7.5) (4.2,8.8) (3.6,10.2) (3.3,11.7) (3.0,13.2) (2.8,14.6) (2.7,16.1) (2.6,17.6) (2.6,19.1) (2.6,20.6) (2.6,22.1) (2.6,23.6) (2.6,25.1) (2.7,26.6) (2.7,28.1) (2.8,29.9)}; \draw[sxb tip] (3.66,9.91) -- (3.56,10.40);
\draw[sxb line] plot[smooth] coordinates {(15.0,0.0) (13.1,-1.1) (11.8,-1.9) (10.5,-2.6) (9.3,-3.5) (8.0,-4.3) (6.9,-5.3) (5.8,-6.3) (4.9,-7.5) (4.2,-8.8) (3.6,-10.2) (3.3,-11.7) (3.0,-13.2) (2.8,-14.6) (2.7,-16.1) (2.6,-17.6) (2.6,-19.1) (2.6,-20.6) (2.6,-22.1) (2.6,-23.6) (2.6,-25.1) (2.7,-26.6) (2.7,-28.1) (2.8,-29.9)}; \draw[sxb tip] (3.66,-9.91) -- (3.56,-10.40);
\draw[sxb line] plot[smooth] coordinates {(15.0,0.0) (13.9,-1.9) (13.2,-3.2) (12.5,-4.5) (11.8,-5.9) (11.2,-7.3) (10.6,-8.7) (10.2,-10.1) (9.8,-11.5) (9.6,-13.0) (9.4,-14.5) (9.3,-16.0) (9.2,-17.5) (9.3,-19.0) (9.3,-20.5) (9.4,-22.0) (9.6,-23.5) (9.7,-25.0) (9.9,-26.5) (10.1,-27.9) (10.5,-29.9)}; \draw[sxb tip] (9.64,-12.70) -- (9.57,-13.20);
\draw[sxb line] plot[smooth] coordinates {(15.0,0.0) (15.0,-2.2) (15.0,-3.7) (15.1,-5.2) (15.1,-6.7) (15.3,-8.2) (15.4,-9.7) (15.6,-11.2) (15.9,-12.6) (16.2,-14.1) (16.6,-15.6) (17.0,-17.0) (17.4,-18.5) (17.8,-19.9) (18.3,-21.3) (18.8,-22.7) (19.4,-24.1) (19.9,-25.5) (20.5,-26.9) (21.1,-28.3) (21.8,-29.9)}; \draw[sxb tip] (16.12,-13.81) -- (16.25,-14.29);
\draw[sxb line] plot[smooth] coordinates {(15.0,0.0) (16.1,-1.9) (16.9,-3.2) (17.6,-4.5) (18.4,-5.8) (19.3,-7.0) (20.1,-8.3) (20.9,-9.5) (21.8,-10.7) (22.7,-11.9) (23.6,-13.1) (24.6,-14.3) (25.5,-15.4) (26.5,-16.6) (27.5,-17.7) (28.4,-18.8) (29.4,-20.0) (30.4,-21.1) (31.5,-22.2) (32.5,-23.3) (33.5,-24.4) (34.5,-25.5) (35.6,-26.5) (36.6,-27.6) (37.7,-28.7) (38.9,-29.9)}; \draw[sxb tip] (24.41,-14.07) -- (24.73,-14.45);
\draw[sxb line] plot[smooth] coordinates {(15.0,0.0) (16.9,-1.1) (18.2,-1.8) (19.5,-2.6) (20.8,-3.3) (22.1,-4.0) (23.5,-4.7) (24.8,-5.4) (26.1,-6.1) (27.5,-6.8) (28.8,-7.5) (30.2,-8.1) (31.5,-8.8) (32.9,-9.4) (34.2,-10.1) (35.6,-10.7) (36.9,-11.3) (38.3,-12.0) (39.7,-12.6) (41.0,-13.2) (42.4,-13.8) (43.8,-14.4) (45.1,-15.0) (46.5,-15.6) (47.9,-16.2) (49.3,-16.8) (50.7,-17.4) (51.8,-17.9)}; \draw[sxb tip] (31.22,-8.68) -- (31.68,-8.88);
  \end{scope}
  \draw[bkGray, line width=0.4pt] (-52,-30) rectangle (52,30);
  \sxbplus{(-15,0)}{3.2}\sxbplus{(15,0)}{3.2}
  \fill[bkInk] (0,0) circle (0.9);
  \node[sxb lab, anchor=north] at (0,-31) {正と正（同じ大きさ）};
\end{scope}
\end{tikzpicture}
\bkfigcaption{2つの点電荷の電気力線（矢じりは電場の向き）。左：正から出た線が負に入る。右：正どうしでは線がはね返し合うように外へ曲がり、まん中の点（黒い点）では電場が0になる。Python で電場の矢印を少しずつつないでかいた（⑩）。}\label{ph:電-1:fig-rikisen}
\end{bkfigure}

\noindent 図\ref{ph:電-1:fig-rikisen}の左では、正の電荷から出た線が、負の電荷に吸いこまれます。2つの電荷のあいだに線がこみ合っていて、電場が強くなっています。右の正どうしでは、線は相手をさけるように外へ曲がり、2つのまん中には線が来ません。まん中では、2つの電荷の電場が同じ大きさで反対向きなので、打ち消し合って0になります。

電気力線を「ゴムひも」のように見ると、力の向きがわかりやすくなります。線は、縮もうとし（正と負が引き合う）、となりの線とは押し合う（正と正がはね返す）、と考えるのです。これは、ファラデーが考えた見方で、電-3 の極板が引き合う力（⑪）や、電場がエネルギーをもつという考え（電-3・電-14）につながります。

\begin{seikaku}
電気力線の「本数」は、実際に数えられる線があるわけではなく、電場の強さを表すための約束です。電-2 では、「電場に垂直な面 \qty{1}{m^2} を通る本数が、電場の強さ $E$ に等しい」と決めて、本数を数えます。すると、電荷 $Q$ から出る本数は $Q/\varepsilon_0=4\pi kQ$ 本になります（高校の教科書の「$4\pi kQ$ 本」）。
\end{seikaku}

%% ---------------------------------------------------------------
\bksubsection{6}{電気双極子：遠くで速く弱まる電場}
\label{ph:電-1:sec6}
大きさが同じで符号が違う2つの電荷 $+q$ と $-q$ が、短い距離 $d$ だけ離れて組になったものを、\textbf{電気双極子}といいます。\bkidx{でんきそうきょくし}{電気双極子}全体の電荷は0ですが、正の電荷と負の電荷の中心がずれています。$-q$ から $+q$ へ向かう、大きさ $p=qd$ の矢印 $\vec p$ を、\textbf{双極子モーメント}といいます（単位 C·m）。

\begin{bkfigure}
\begin{tikzpicture}[x=1mm, y=1mm]
  \draw[bkGray, line width=0.4pt, dashed] (-8,0) -- (96,0);
  \sxbminus{(0,0)}{2.4}\sxbplus{(14,0)}{2.4}
  \draw[bkGreen!50!black, line width=1.2pt, -{Stealth[length=2.6mm]}] (2.8,6) -- (11.2,6);
  \node[sxb lab, anchor=south, text=bkGreen!40!black] at (7,6.6) {$\vec p$};
  \draw[sxb dim] (0,-9) -- (14,-9);
  \node[sxb lab, anchor=north] at (7,-9.5) {$d$};
  \fill[bkInk] (7,0) circle (0.6);
  \fill[bkInk] (80,0) circle (1);
  \draw[sxb dim] (7,-16) -- (80,-16);
  \node[sxb lab, anchor=north] at (43.5,-16.5) {$r$（まん中から）};
  \draw[bkRed!80!black, sxb f] (80,4) -- (92,4);
  \node[sxb lab, anchor=south, text=bkRed!75!black] at (86,4.6) {$+q$ から（強い）};
  \draw[bkBlue!80!black, sxb f] (80,-4) -- (70,-4);
  \node[sxb lab, anchor=north, text=bkBlue!75!black] at (72,-4.6) {$-q$ へ（少し弱い）};
  \node[sxb lab, anchor=east] at (-2.8,-3) {$-q$};
  \node[sxb lab, anchor=west] at (16.8,-3) {$+q$};
\end{tikzpicture}
\bkfigcaption{電気双極子の軸の上の電場。$+q$ のほうが少し近いので、$+q$ からの外向きの電場が、$-q$ へ向かう電場に少しだけ勝つ。その「少し」が、遠くほど小さくなる。}\label{ph:電-1:fig-soukyoku}
\end{bkfigure}

\phboxhead{軸の上の電場}
双極子のまん中から、軸の上に距離 $r$ 離れた点（$+q$ の側）を考えます（図\ref{ph:電-1:fig-soukyoku}）。$+q$ までの距離は $r-d/2$、$-q$ までは $r+d/2$ で、2つの電場は軸の上で反対向きです。重ね合わせで、
\[E=\frac{kq}{(r-d/2)^2}-\frac{kq}{(r+d/2)^2}\]
です。$r$ が $d$ よりずっと大きい遠くでは、5-17 の近似 $(1+x)^{-2}\approx1-2x$ が使えます。
\[\frac{1}{(r\mp d/2)^2}=\frac{1}{r^2}\Bigl(1\mp\frac{d}{2r}\Bigr)^{-2}\approx\frac{1}{r^2}\Bigl(1\pm\frac{d}{r}\Bigr)\]
だから、2つの差は $\dfrac{kq}{r^2}\times\dfrac{2d}{r}$ です。
\pheq{E\approx\frac{2kqd}{r^3}=\frac{2kp}{r^3}}{N/C ＝ $\dfrac{\text{N·m}^2}{\text{C}^2}$×$\dfrac{\text{C·m}}{\text{m}^3}$}
\label{ph:電-1:dipole}電場は、距離の\textbf{3乗}に反比例します。


\begin{phyomi}[$E\approx\dfrac{2kp}{r^3}$]
「遠くから見ると、正と負がほとんど重なって打ち消し合うので、電場は点電荷より速く弱まる」。1つの電荷の電場（$1/r^2$）に、「どれだけずれているか」の割合 $d/r$ がかかった形です。
\end{phyomi}

\begin{phbaisuru}{距離 $r$}
電場は $\tfrac18$ になります。点電荷の $\tfrac14$ より、ずっと速く弱まります（例題8）。下の表は、±1\,nC・$d=\qty{1}{cm}$ の双極子の軸の上の電場です。近似の式は、\qty{0.1}{m}でも0.5\,\% しかずれていません。
\par\smallskip\noindent\begin{bkcenter}
\begin{bktable}{colspec={Q[c,m,wd=34mm] Q[c,m,wd=26mm] Q[c,m,wd=26mm]}, rowsep=1.5pt, row{1}={bg=bkTheme!10, font=\sffamily\small}}
距離 $r$〔m〕 & 0.1 & 0.2 \\
重ね合わせ（正確）〔N/C〕 & 181 & 22.5 \\
近似 $2kp/r^3$〔N/C〕 & 180 & 22.5 \\
\end{bktable}
\end{bkcenter}
\end{phbaisuru}

\phboxhead{水は双極子}
水の分子（H$_2$O）は、酸素の側が少し負、水素の側が少し正になっていて、1つ1つが小さな双極子です。やってみよう①で、水の細い流れが下じきへ曲がったのは、このためです。下じきの近くでは、水の分子が、正の側（または負の側）を下じきへ向けるように回ります。すると、下じきに近い側の電荷のほうが強く引かれ、遠い側の電荷がはね返される力より大きくなるので、全体として引き寄せられるのです。\textbf{電場が場所によって違う}（近いほど強い）ことが大切です。一様な電場の中では、双極子は回るだけで、引き寄せられません。電子レンジは、1秒に約24億5千万回、向きが行ったり来たりする電場で水の分子をゆすって、温めています（電-15）。ちぎったティッシュが引きつけられるのも、ティッシュの中の電荷が少しずれて、双極子のようになるからです（電-3 の静電誘導・誘電分極）。

\begin{motto}[4]
\phboxhead{一様な電場の中の双極子：回るけれど、進まない}
\noindent\begin{minipage}[c]{0.6\linewidth}\raggedright
一様な電場 $E$ の中に、双極子（$\pm q$、間かく $d$）を、電場から角 $\theta$ だけ傾けて置きます。$+q$ は電場の向きに $qE$、$-q$ は反対向きに $qE$ の力を受けます。2つの力は大きさが同じで向きが反対なので、合わせた力は0で、双極子は動き出しません。でも、2つの力の作用線がずれているので、双極子を回そうとします。回そうとする力のモーメント（トルク：力-6）の大きさは、力 $qE$ に腕の長さ $d\sin\theta$ をかけて、
\end{minipage}\hfill
\begin{minipage}[c]{0.37\linewidth}\centering
\begin{tikzpicture}[x=1mm, y=1mm]
  \foreach \yy in {-5,23} {\draw[sxb arr] (-2,\yy) -- (40,\yy);}
  \node[sxb lab, anchor=west, text=bkNavy] at (40.5,23) {$\vec E$};
  \draw[bkGray, line width=0.4pt, dashed] (12,4) -- (32,4);
  \draw[bkInk, line width=0.8pt] (12,4) -- (28,14);
  \draw[bkInk, line width=0.5pt] (20,4) arc (0:32:8);
  \node[sxb lab] at (22.6,6.4) {$\theta$};
  \sxbminus{(12,4)}{2.2}\sxbplus{(28,14)}{2.2}
  \draw[bkRed!80!black, sxb f] (30.4,14) -- (38,14);
  \draw[bkRed!80!black, sxb f] (9.6,4) -- (2,4);
\end{tikzpicture}
\end{minipage}
\[qE\times d\sin\theta=pE\sin\theta\]
です。このモーメントは、$\vec p$ を電場の向きにそろえようとします（$\theta=0$ で0）。電子レンジの中で、向きがすばやく入れかわる電場に合わせて水の分子が回ろうとし、まわりの分子とこすれて熱が出るのは、このモーメントのためです（電-15）。一様でない電場の中では、力の大きさが $+q$ と $-q$ で少し違うので、回るだけでなく、電場の強い所へ引き寄せられます（やってみようの水の流れ）。
\end{motto}

%% ---------------------------------------------------------------
\Needspace{12\baselineskip}\bksubsection{7}{広がった電荷の電場：積分で重ね合わせる}
\label{ph:電-1:sec7}
これまでは、点電荷だけを考えました。でも、帯電した下じきや、電気を帯びた長い電線のように、電荷が広がっている物もあります。そのときは、電荷を小さく分けて、1つ1つを点電荷とみなし、それぞれの電場を重ね合わせます。分け方を細かくした極限が、積分（8-6）です。電荷の広がり方に合わせて、次の密度を使います。
\begin{itemize}[leftmargin=1.2\zw, itemsep=1pt, topsep=2pt, label={\color{bkTheme}●}]
\item 線に広がった電荷：1 m あたりの電荷 $\lambda$（ラムダ）〔C/m〕。\textbf{線電荷密度}。
\item 面に広がった電荷：1 m$^2$ あたりの電荷 $\sigma$（シグマ）〔C/m$^2$〕。\textbf{面電荷密度}。
\end{itemize}
\bkidx{せんでんかみつど}{線電荷密度}\bkidx{めんでんかみつど}{面電荷密度}

\phboxhead{まっすぐな長い線}
\noindent\begin{minipage}[c]{0.5\linewidth}\raggedright
まっすぐな線（$x$ 軸）に $\lambda$ が一様に広がっています。線から $r$ 離れた点 P の電場を求めます。線の上の $x$ の所の短い長さ $\dd x$ には電荷 $\lambda\,\dd x$ があり、P までの距離は $\sqrt{x^2+r^2}$ です。その電場の、線に垂直な成分は、$\cos\theta=r/\sqrt{x^2+r^2}$ をかけて
\end{minipage}\hfill
\begin{minipage}[c]{0.47\linewidth}\centering
\begin{tikzpicture}[x=1mm, y=1mm]
  \draw[bkOrange!80!black, line width=2pt] (-26,0) -- (26,0);
  \foreach \xx in {-22,-15,...,22} {\sxbp{(\xx,-2.6)}}
  \fill[bkInk] (0,22) circle (0.9);
  \node[sxb lab, anchor=south] at (0,23) {P};
  \draw[sxb dim] (-3,0) -- (-3,22);
  \node[sxb lab, anchor=east] at (-3.5,11) {$r$};
  \fill[bkRed!70!black] (13,-0.8) rectangle (17,0.8);
  \node[sxb lab, anchor=north] at (15,-3.6) {$\lambda\,\dd x$};
  \draw[bkGray, line width=0.4pt, dashed] (15,0) -- (0,22);
  \draw[bkRed!80!black, sxb f] (0,22) -- ++({-15/26.6*9},{22/26.6*9});
  \draw[bkInk, line width=0.4pt] (0,16) arc (-90:-55.7:6);
  \node[sxb lab] at (3,13.6) {$\theta$};
  \draw[sxb dim] (0,-8.5) -- (15,-8.5);
  \node[sxb lab, anchor=north] at (7.5,-9) {$x$};
\end{tikzpicture}
\bkfigcaption{長い線の電荷がつくる電場。}\label{ph:電-1:fig-sen}
\end{minipage}
\[\dd E_\perp=k\,\frac{\lambda\,\dd x}{x^2+r^2}\cdot\frac{r}{\sqrt{x^2+r^2}}=k\lambda r\,\frac{\dd x}{(x^2+r^2)^{3/2}}\]
です。線に平行な成分は、P の左右の対称な2か所で打ち消し合うので、垂直な成分だけをたせばよいのです。積分の公式 $\displaystyle\int_{-\infty}^{\infty}\frac{\dd x}{(x^2+r^2)^{3/2}}=\frac{2}{r^2}$（$x=r\tan\theta$ と置きかえる：5-11）を使うと、
\pheq{E=k\lambda r\times\frac{2}{r^2}=\frac{2k\lambda}{r}=\frac{\lambda}{2\pi\varepsilon_0r}}{N/C ＝ $\dfrac{\text{N·m}^2}{\text{C}^2}$×$\dfrac{\text{C/m}}{\text{m}}$}
\label{ph:電-1:sen}です。長い線の電場は、距離に（2乗ではなく）\textbf{1乗}で反比例します。長さが $L$ の線なら、積分の範囲を $-L/2$ から $L/2$ にして、
\[E=\frac{2k\lambda}{r}\times\frac{L/2}{\sqrt{r^2+(L/2)^2}}\]
です。$L$ が $r$ よりずっと長ければ、後ろの分数はほぼ1になります（章末問題）。

\phboxhead{円環の軸の上}
半径 $a$ の細い輪（円環）に、電荷 $Q$ が一様に広がっています。輪の中心から軸の上に $z$ 離れた点では、輪のどの部分からの距離も $\sqrt{z^2+a^2}$ で同じです。軸に垂直な成分は、輪の反対側の部分と打ち消し合い、軸の向きの成分（$\cos$ の割合 $z/\sqrt{z^2+a^2}$）だけが残ります。どの部分も同じ割合なので、積分は電荷 $Q$ をたすだけになり、
\pheq{E=\frac{kQz}{(z^2+a^2)^{3/2}}}{N/C}
\label{ph:電-1:enkan}です。$z=0$（輪の中心）では0、$z$ が $a$ よりずっと大きい遠くでは $kQ/z^2$ で、点電荷と同じになります。0から始まって、いったん大きくなり、また小さくなるので、どこかでいちばん強くなります。$E$ を $z$ で微分して0とおくと（5-7）、
\[\frac{\dd E}{\dd z}=kQ\,\frac{(z^2+a^2)-3z^2}{(z^2+a^2)^{5/2}}=0\qquad\text{より}\qquad z=\frac{a}{\sqrt2}\]
です（例題10）。

\begin{bkfigure}
\begin{tikzpicture}
\begin{axis}[bkplot, width=112mm, height=46mm, xmin=0, xmax=0.42, ymin=0, ymax=4200, xtick={0.0707,0.1,0.2,0.3,0.4}, xticklabels={$\tfrac{a}{\sqrt2}$,0.1,0.2,0.3,0.4}, ytick={1000,2000,3460}, xlabel={$z$（m）}, ylabel={$E$（N/C）}]
  \addplot[bkcurve, domain=0:0.4, samples=120] {89.9*x/(x^2+0.01)^1.5};
  \addplot[bkOrange, line width=1.4pt, dash pattern=on 4pt off 2pt, domain=0.18:0.4, samples=60] {89.9/x^2};
  \fill[bkRed] (axis cs:0.0707,3460) circle (1.8pt);
  \node[font=\sffamily\footnotesize, anchor=north west, fill=white, inner sep=1pt] at (axis cs:0.085,4150) {いちばん強い：$3.5\times10^3$};
  \node[font=\sffamily\footnotesize, anchor=west, fill=white, inner sep=1pt] at (axis cs:0.24,3100) {点電荷とみた $kQ/z^2$（破線）};
\end{axis}
\end{tikzpicture}
\bkfigcaption{円環（$Q=\qty{1e-8}{C}$、半径 $a=\qty{0.1}{m}$）の軸の上の電場（実線）。中心では0、$z=a/\sqrt2\approx\qty{0.071}{m}$ でいちばん強い。遠くでは、点電荷の電場（破線）に近づく。}\label{ph:電-1:fig-enkan}
\end{bkfigure}

\phboxhead{広い平らな面：距離によらない電場}
面電荷密度 $\sigma$ の円板（半径 $R$）を、半径 $s$・はば $\dd s$ の細い円環に分けます。1つの円環の電荷は $\sigma\times2\pi s\,\dd s$ です。円環の式に入れて、$s$ を0から $R$ まで積分すると、軸の上で
\[E=\int_0^R\frac{kz\,\sigma\,2\pi s\,\dd s}{(z^2+s^2)^{3/2}}=2\pi k\sigma\Bigl(1-\frac{z}{\sqrt{z^2+R^2}}\Bigr)=\frac{\sigma}{2\varepsilon_0}\Bigl(1-\frac{z}{\sqrt{z^2+R^2}}\Bigr)\]
です（$2\pi k=1/(2\varepsilon_0)$）。$R$ を限りなく大きくした「どこまでも広い平らな面」では、かっこの中が1になり、
\pheq{E=\frac{\sigma}{2\varepsilon_0}}{N/C ＝ $\dfrac{\text{C/m}^2}{\text{F/m}}$}
\label{ph:電-1:heimen}です。\textbf{面からの距離 $z$ が入っていません}。広い面のすぐ近くでも、少し離れても、同じ強さの電場なのです。面の両側に、面に垂直に外向き（$\sigma>0$ のとき）の電場ができます。

\begin{phflow}[広がった電荷の電場]
\phflowrow{$\vec E=\displaystyle\int k\,\frac{\dd Q}{r^2}\,\vec e_r$}{まっすぐな\\長い線}{$E=\dfrac{\lambda}{2\pi\varepsilon_0r}$}
\phflowrow{$\vec E=\displaystyle\int k\,\frac{\dd Q}{r^2}\,\vec e_r$}{どこまでも広い\\平らな面}{$E=\dfrac{\sigma}{2\varepsilon_0}$}
\phflowrow{$E=\dfrac{\sigma}{2\varepsilon_0}$（2枚）}{正と負の面を\\向かい合わせる}{間は $E=\dfrac{\sigma}{\varepsilon_0}$、外は0}
\end{phflow}

\begin{bkfigure}
\begin{tikzpicture}[x=1mm, y=1mm]
  \fill[bkRed!12] (19,0) rectangle (23,30);
  \fill[bkBlue!12] (55,0) rectangle (59,30);
  \draw[bkRed!70!black, line width=0.6pt] (19,0) rectangle (23,30);
  \draw[bkBlue!70!black, line width=0.6pt] (55,0) rectangle (59,30);
  \foreach \yy in {3,9,...,27} {\sxbp{(21,\yy)}\sxbm{(57,\yy)}}
  \foreach \yy in {5,15,25} {
    \draw[bkRed!80!black, line width=0.8pt, -{Stealth[length=2mm]}] (24,{\yy+1.2}) -- (37,{\yy+1.2});
    \draw[bkRed!80!black, line width=0.8pt, -{Stealth[length=2mm]}] (18,{\yy+1.2}) -- (5,{\yy+1.2});
    \draw[bkBlue!80!black, line width=0.8pt, dash pattern=on 2pt off 1pt, -{Stealth[length=2mm]}] (41,{\yy-1.2}) -- (54,{\yy-1.2});
    \draw[bkBlue!80!black, line width=0.8pt, dash pattern=on 2pt off 1pt, -{Stealth[length=2mm]}] (73,{\yy-1.2}) -- (60,{\yy-1.2});}
  \node[sxb lab, anchor=north] at (21,-1) {$+\sigma$};
  \node[sxb lab, anchor=north] at (57,-1) {$-\sigma$};
  \node[anchor=west, font=\small, text width=46mm, align=left] at (78,15) {%
  実線：正の面の電場（両側に外向き）。破線：負の面の電場（両側から面へ向かう）。\par\vspace{2pt}
  あいだでは同じ向きで2倍、外側では反対向きで打ち消し合う。};
\end{tikzpicture}
\bkfigcaption{向かい合った正と負の広い面。あいだの電場は $\sigma/(2\varepsilon_0)\times2=\sigma/\varepsilon_0$ で一様、外側は0。電-3 のコンデンサーの中の電場は、このしくみでできる。}\label{ph:電-1:fig-2mai}
\end{bkfigure}

\begin{phbaisuru}[半分]{円板の半径 $R$}
平らな面の式は「どこまでも広い」ときの式でした。本物の板は有限です。半径 $R=\qty{0.1}{m}$ の円板では、軸の上の電場は、広い面の値に対して、$z=\qty{1}{mm}$ で99\,\%、\qty{1}{cm}で90\,\%、\qty{5}{cm}で55\,\%、\qty{10}{cm}で29\,\% です。半径を半分の\qty{0.05}{m}にすると、同じ割合になる距離も半分になります。\textbf{板の大きさよりずっと近い所}では、「広い面」とみなしてよいのです。
\end{phbaisuru}

\noindent 線・面の電場は、電-2 のガウスの法則を使うと、積分をほとんどしないで同じ答えが出ます（電-2 ④）。ここでの積分は「重ね合わせの原理から、まっすぐに求めると、こうなる」という確かめです。どちらで求めても同じ答えになることが、2つの法則が同じ内容であることを示しています。

%% ---------------------------------------------------------------
\Needspace{12\baselineskip}\bksubsection{8}{電場の中の電荷の動き}
\label{ph:電-1:sec8}
向かい合った広い2枚の板のあいだのように、電場の大きさと向きがどこでも同じ所を、\textbf{一様な電場}といいます。\bkidx{いちようなでんば}{一様な電場}一様な電場 $\vec E$ の中で、質量 $m$・電荷 $q$ の粒は、一定の力 $q\vec E$ を受けます。運動方程式（力-2）は
\pheq{m\vec a=q\vec E\qquad\text{だから}\qquad\vec a=\frac{q}{m}\vec E}{m/s² ＝ $\dfrac{\text{C}}{\text{kg}}$×$\dfrac{\text{N}}{\text{C}}$}
で、加速度は一定です。これは、重力（$m\vec a=m\vec g$）の中の物の動きと同じ形です。電場に垂直に飛びこんだ粒は、力-1 ④の投げた物と同じように、\textbf{放物線}をえがいて曲がります。違うのは、加速度の大きさが $q/m$（\textbf{比電荷}）で決まり、粒の重さによって変わることです。電子は、比電荷がとても大きい（$e/m_{\mathrm e}\approx1.76\times10^{11}$\,C/kg）ので、弱い電場でも大きく曲がります。

\begin{phflow}[一様な電場の中の粒]
\phflowrow{$m\vec a=q\vec E$}{電場が一様\\重力は小さい}{$a=\dfrac{qE}{m}$（一定）}
\phflowrow{$a=\dfrac{qE}{m}$（一定）}{電場に垂直に\\速さ $v_0$ で入る}{$y=\dfrac12\cdot\dfrac{qE}{m}\Bigl(\dfrac{x}{v_0}\Bigr)^2$}
\end{phflow}

\begin{bkfigure}
\begin{tikzpicture}[x=1mm, y=1mm]
  \fill[bkBlue!20] (30,22) rectangle (80,24);
  \fill[bkRed!20] (30,-2) rectangle (80,0);
  \draw[bkBlue!70!black, line width=0.6pt] (30,22) rectangle (80,24);
  \draw[bkRed!70!black, line width=0.6pt] (30,-2) rectangle (80,0);
  \foreach \xx in {34,42,...,78} {\sxbm{(\xx,25.8)}\sxbp{(\xx,-3.8)}}
  \foreach \xx in {38,52,66} {\draw[sxb arr] (\xx,1.5) -- (\xx,7);}
  \node[sxb lab, anchor=west, text=bkNavy] at (67,4.5) {$\vec E$};
  % 滴の道
  \draw[bkGray, line width=0.5pt, dashed] (0,14) -- (124,14);
  \draw[bkInk, line width=1.1pt] (0,14) -- (30,14);
  \draw[bkInk, line width=1.1pt, domain=0:50, samples=40] plot ({30+\x}, {14-0.0024*\x*\x});
  \draw[bkInk, line width=1.1pt, -{Stealth[length=2.4mm]}] (80,8) -- (110,0.8);
  \filldraw[fill=bkBlue!60, draw=bkInk, line width=0.4pt] (8,14) circle (1.5);
  \node[sxb lab, anchor=north] at (8,11.6) {負に帯電した滴};
  \draw[bkRed!80!black, line width=0.8pt, -{Stealth[length=2mm]}] (13,17.5) -- (23,17.5);
  \node[sxb lab, anchor=south, text=bkRed!75!black] at (18,18) {$v_0$};
  \node[sxb lab, anchor=north] at (55,-5.6) {長さ $l$ の板のあいだ：放物線};
  \node[sxb lab, anchor=north west] at (86,-0.5) {出たあとはまっすぐ};
\end{tikzpicture}
\bkfigcaption{一様な電場の中を通る、帯電した粒。負に帯電した粒は、電場（上向き）と反対の下向きに一定の加速度を受けて、放物線をえがく。インクジェットプリンターや、昔のテレビ（ブラウン管）は、このしくみで粒の行き先を決めた。}\label{ph:電-1:fig-ink}
\end{bkfigure}

\noindent 板のあいだを通る時間は $t=l/v_0$ で、そのあいだに、電場の向きに $\tfrac12at^2$ ずれます。電場を強くすると、ずれは大きくなります。電場の向きを変えれば、反対側へ曲げられます。ある種のインクジェットプリンターは、小さなインクの滴に、印刷したい点に合わせて違う量の電荷をのせ、板のあいだで曲げる量を変えて、紙の上の位置を決めます（例題12）。

\begin{phmodel}[一様な電場の中の粒]
\phignore{粒にはたらく重力（例題12 で、電気の力の約200分の1）。空気の抵抗。板のはしで電場が弱まる所。}
\phassume{板のあいだの電場は一様で、外は0。粒の電荷と質量は変わらない。}
\phlimit{粒がとても遅く、重いとき（重力が効く）。電子が光の速さに近いとき（相対論：原-11）。}
\end{phmodel}

%% ---------------------------------------------------------------
\bksubsection{9}{静電気を使う・静電気から守る}
\label{ph:電-1:sec9}
静電気の力は、身の回りの機械で、上手に使われています。
\begin{itemize}[leftmargin=1.2\zw, itemsep=2pt, topsep=2pt, label={\color{bkTheme}●}]
\item \textbf{複写機・レーザープリンター}：光を当てた所だけ電荷が消える筒（感光体）に、帯電した細かい粉（トナー）をつけ、紙に移して熱でとかし、固めます。
\item \textbf{静電塗装}：塗料の霧を帯電させ、反対の符号にした車体へ引きつけます。塗料が裏側にも回りこみ、むだが少なくなります。
\item \textbf{電気集じん器・空気清浄機}：空気の中のちりや煙の粒を帯電させ、反対の符号の板に集めます。発電所の煙突の手前でも使われています。
\end{itemize}

\noindent いっぽう、静電気は困ったことも起こします。冬にドアノブに触れたときの「パチッ」は、歩くあいだに帯電した体から、電荷が一気に飛ぶ\textbf{放電}です。人がはっきり感じるのは、体の電圧（電位：電-2）がおよそ\qty{3}{kV}（3000 V）をこえたときといわれます。小麦粉や石炭の粉が舞う工場では、放電の火花で\textbf{粉じん爆発}が起きることがあります。セルフ式のガソリンスタンドには、給油の前に触れて体の電荷を地面へ逃がす\textbf{静電気除去シート}があります。燃えやすいガソリンの蒸気に、放電の火花が飛ばないようにするためです。スマホやパソコンの中の小さな部品も、静電気の放電でこわれることがあるので、工場では、手首に地面とつないだバンドをつけて作業します。

%% ---------------------------------------------------------------
\Needspace{14\baselineskip}\bksubsection{10}{表計算と Python で、電場の地図をかく}
\label{ph:電-1:sec10}
いくつかの電荷がつくる電場を、場所ごとに計算して、地図にしましょう。8-5⑤ と同じ方法です。$+1$\,nC を $(-0.1,\ 0)$、$-1$\,nC を $(0.1,\ 0)$〔m〕に置きます。場所 $(x,\ y)$ での電場の成分は、2つの点電荷の電場を成分に分けてたして、
\[E_x=kq_1\frac{x-a_1}{r_1^{\,3}}+kq_2\frac{x-a_2}{r_2^{\,3}},\qquad E_y=kq_1\frac{y}{r_1^{\,3}}+kq_2\frac{y}{r_2^{\,3}}\]
です（$a_1=-0.1$、$a_2=0.1$、$r_1$・$r_2$ は、それぞれの電荷からの距離）。$\dfrac{x-a}{r}$ が向きの $\cos$、$\dfrac{1}{r^2}$ が大きさの分なので、合わせて $\dfrac{x-a}{r^3}$ になります。

表計算では、A列に $x$、B列に $y$ を入れ、C列・D列に $r_1$・$r_2$、E列・F列に $E_x$・$E_y$ の式を入れます。
\begin{bkcenter}\small
\texttt{C2 = SQRT((A2+0.1)\textasciicircum2+B2\textasciicircum2)}\quad\texttt{D2 = SQRT((A2-0.1)\textasciicircum2+B2\textasciicircum2)}\\
\texttt{E2 = 8.99*((A2+0.1)/C2\textasciicircum3-(A2-0.1)/D2\textasciicircum3)}\\
\texttt{F2 = 8.99*(B2/C2\textasciicircum3-B2/D2\textasciicircum3)}
\end{bkcenter}
$k\times\qty{1e-9}{C}=8.99$ なので、式の前に8.99をかけています。

\begin{bkcenter}
\bksheet{6}{15}{%
1 & $x$〔m〕 & $y$〔m〕 & $r_1$〔m〕 & $r_2$〔m〕 & $E_x$〔N/C〕 & $E_y$〔N/C〕 \\
2 & 0 & 0 & 0.1 & 0.1 & 1798 & 0 \\
3 & 0 & 0.1 & 0.1414 & 0.1414 & 636 & 0 \\
4 & 0.1 & 0.1 & 0.2236 & 0.1 & 161 & $-819$ \\
5 & $-0.2$ & 0 & 0.1 & 0.3 & $-799$ & 0 \\
6 & 0 & 0.2 & 0.2236 & 0.2236 & 161 & 0 \\}
\bktabcaption{電場の表計算（2つの点電荷 $\pm1$\,nC）。2行め（まん中）では、2つの電場が同じ向き（右向き）でたし合わされて強い。5行め（正の電荷の左）では、左向き。}\label{ph:電-1:tab-sheet}
\end{bkcenter}

\noindent 2行めを手で確かめます。まん中 $(0,\ 0)$ は、どちらの電荷からも\qty{0.1}{m}です。$+1$\,nC の電場は右向きに899、$-1$\,nC の電場は（負の電荷へ向かうので）これも右向きに899で、合わせて右向きに\qty{1798}{N/C}です。

\noindent Python では、同じ計算を関数にして、たくさんの場所で繰り返します。
\begin{phcode}[Python]
import math
k = 8.99e9
qs = [(1e-9, -0.1, 0.0), (-1e-9, 0.1, 0.0)]   \# (電荷, x, y)
def E(x, y):
    Ex = Ey = 0.0
    for q, a, b in qs:                         \# 電荷ごとにたす
        r = math.hypot(x - a, y - b)
        Ex += k * q * (x - a) / r**3
        Ey += k * q * (y - b) / r**3
    return Ex, Ey
for x, y in [(0, 0), (0, 0.1), (0.1, 0.1)]:
    Ex, Ey = E(x, y)
    print(x, y, round(Ex), round(Ey), round(math.hypot(Ex, Ey)))
\end{phcode}
\begin{phcodeout}
0 0 1798 0 1798
0 0.1 636 0 636
0.1 0.1 161 -819 834
\end{phcodeout}
左から、$x$、$y$、$E_x$、$E_y$、電場の大きさ〔N/C〕です。表計算と同じ数になりました。\texttt{qs} の行に電荷を書き足せば、3つでも10個でも、同じプログラムで計算できます。

\begin{bkfigure}
\begin{tikzpicture}[x=1mm, y=1mm]
  \draw[bkGray, line width=0.4pt] (-55,-34) rectangle (55,34);
\draw[sxb arr] (-49.3,-29.0) -- (-50.7,-31.0);
\draw[sxb arr] (-48.9,-19.1) -- (-51.1,-20.9);
\draw[sxb arr] (-48.5,-9.4) -- (-51.5,-10.6);
\draw[sxb arr] (-48.4,0.0) -- (-51.6,0.0);
\draw[sxb arr] (-48.5,9.4) -- (-51.5,10.6);
\draw[sxb arr] (-48.9,19.1) -- (-51.1,20.9);
\draw[sxb arr] (-49.3,29.0) -- (-50.7,31.0);
\draw[sxb arr] (-39.5,-28.7) -- (-40.5,-31.3);
\draw[sxb arr] (-38.9,-18.7) -- (-41.1,-21.3);
\draw[sxb arr] (-38.2,-9.0) -- (-41.8,-11.0);
\draw[sxb arr] (-37.8,0.0) -- (-42.2,0.0);
\draw[sxb arr] (-38.2,9.0) -- (-41.8,11.0);
\draw[sxb arr] (-38.9,18.7) -- (-41.1,21.3);
\draw[sxb arr] (-39.5,28.7) -- (-40.5,31.3);
\draw[sxb arr] (-29.9,-28.4) -- (-30.1,-31.6);
\draw[sxb arr] (-29.3,-18.1) -- (-30.7,-21.9);
\draw[sxb arr] (-28.1,-7.9) -- (-31.9,-12.1);
\draw[sxb arr] (-26.4,0.0) -- (-33.6,0.0);
\draw[sxb arr] (-28.1,7.9) -- (-31.9,12.1);
\draw[sxb arr] (-29.3,18.1) -- (-30.7,21.9);
\draw[sxb arr] (-29.9,28.4) -- (-30.1,31.6);
\draw[sxb arr] (-20.6,-28.5) -- (-19.4,-31.5);
\draw[sxb arr] (-20.4,-17.8) -- (-19.6,-22.2);
\draw[sxb arr] (-20.2,-6.4) -- (-19.8,-13.6);
\draw[sxb arr] (-20.2,6.4) -- (-19.8,13.6);
\draw[sxb arr] (-20.4,17.8) -- (-19.6,22.2);
\draw[sxb arr] (-20.6,28.5) -- (-19.4,31.5);
\draw[sxb arr] (-11.3,-29.0) -- (-8.7,-31.0);
\draw[sxb arr] (-11.6,-18.6) -- (-8.4,-21.4);
\draw[sxb arr] (-12.4,-8.3) -- (-7.6,-11.7);
\draw[sxb arr] (-13.8,0.0) -- (-6.2,0.0);
\draw[sxb arr] (-12.4,8.3) -- (-7.6,11.7);
\draw[sxb arr] (-11.6,18.6) -- (-8.4,21.4);
\draw[sxb arr] (-11.3,29.0) -- (-8.7,31.0);
\draw[sxb arr] (-1.6,-30.0) -- (1.6,-30.0);
\draw[sxb arr] (-2.0,-20.0) -- (2.0,-20.0);
\draw[sxb arr] (-2.6,-10.0) -- (2.6,-10.0);
\draw[sxb arr] (-2.9,0.0) -- (2.9,0.0);
\draw[sxb arr] (-2.6,10.0) -- (2.6,10.0);
\draw[sxb arr] (-2.0,20.0) -- (2.0,20.0);
\draw[sxb arr] (-1.6,30.0) -- (1.6,30.0);
\draw[sxb arr] (8.7,-31.0) -- (11.3,-29.0);
\draw[sxb arr] (8.4,-21.4) -- (11.6,-18.6);
\draw[sxb arr] (7.6,-11.7) -- (12.4,-8.3);
\draw[sxb arr] (6.2,0.0) -- (13.8,0.0);
\draw[sxb arr] (7.6,11.7) -- (12.4,8.3);
\draw[sxb arr] (8.4,21.4) -- (11.6,18.6);
\draw[sxb arr] (8.7,31.0) -- (11.3,29.0);
\draw[sxb arr] (19.4,-31.5) -- (20.6,-28.5);
\draw[sxb arr] (19.6,-22.2) -- (20.4,-17.8);
\draw[sxb arr] (19.8,-13.6) -- (20.2,-6.4);
\draw[sxb arr] (19.8,13.6) -- (20.2,6.4);
\draw[sxb arr] (19.6,22.2) -- (20.4,17.8);
\draw[sxb arr] (19.4,31.5) -- (20.6,28.5);
\draw[sxb arr] (30.1,-31.6) -- (29.9,-28.4);
\draw[sxb arr] (30.7,-21.9) -- (29.3,-18.1);
\draw[sxb arr] (31.9,-12.1) -- (28.1,-7.9);
\draw[sxb arr] (33.6,0.0) -- (26.4,0.0);
\draw[sxb arr] (31.9,12.1) -- (28.1,7.9);
\draw[sxb arr] (30.7,21.9) -- (29.3,18.1);
\draw[sxb arr] (30.1,31.6) -- (29.9,28.4);
\draw[sxb arr] (40.5,-31.3) -- (39.5,-28.7);
\draw[sxb arr] (41.1,-21.3) -- (38.9,-18.7);
\draw[sxb arr] (41.8,-11.0) -- (38.2,-9.0);
\draw[sxb arr] (42.2,0.0) -- (37.8,0.0);
\draw[sxb arr] (41.8,11.0) -- (38.2,9.0);
\draw[sxb arr] (41.1,21.3) -- (38.9,18.7);
\draw[sxb arr] (40.5,31.3) -- (39.5,28.7);
\draw[sxb arr] (50.7,-31.0) -- (49.3,-29.0);
\draw[sxb arr] (51.1,-20.9) -- (48.9,-19.1);
\draw[sxb arr] (51.5,-10.6) -- (48.5,-9.4);
\draw[sxb arr] (51.6,0.0) -- (48.4,0.0);
\draw[sxb arr] (51.5,10.6) -- (48.5,9.4);
\draw[sxb arr] (51.1,20.9) -- (48.9,19.1);
\draw[sxb arr] (50.7,31.0) -- (49.3,29.0);
  \sxbplus{(-20,0)}{2.6}\sxbminus{(20,0)}{2.6}
  \foreach \xx in {-50,-30,-10,10,30,50} {\pgfmathsetmacro\lab{\xx/200}\node[sxb lab, anchor=north] at (\xx,-34.6) {\pgfmathprintnumber[fixed, precision=2]{\lab}};}
  \node[sxb lab, anchor=west] at (56,-34) {$x$（m）};
  \foreach \yy in {-30,-10,10,30} {\pgfmathsetmacro\lab{\yy/200}\node[sxb lab, anchor=east] at (-55.6,\yy) {\pgfmathprintnumber[fixed, precision=2]{\lab}};}
  \node[sxb lab, anchor=south] at (-55,34.6) {$y$（m）};
\end{tikzpicture}
\bkfigcaption{Python で計算した電場の地図（\qty{0.05}{m}ごとの点）。矢印の向きが電場の向き。矢印の長さは、強さそのものではなく、強さの3乗根に比例させた（そうしないと、遠くの矢印が見えなくなるため）。電荷のすぐ近くは、矢印が長くなりすぎるので、長さの上限を決めた。}\label{ph:電-1:fig-map}
\end{bkfigure}

\noindent 矢印を少しずつつないでいくと、⑤の図\ref{ph:電-1:fig-rikisen}の電気力線がかけます。ある点の電場の向きに、ほんの少し（\qty{0.25}{mm}ぶん）進み、その点でまた電場の向きを計算して進む、を繰り返します（力-2 ⑨のルンゲ・クッタ法で進めました）。

\begin{phsaki}{電場の正体と、物の中の電気の力}
この章では、電場を「1\,C あたりの力」として定義しました。電場は、電荷が止まっていないと、磁場と組になって「電磁場」としてふるまい、エネルギーや運動量を運びます（電-14・電-15）。さらに、量子力学と相対論を合わせた「場の量子論」では、電気の力は、光の粒（光子）をやりとりすることで伝わる、と考えます。また、原子や分子のあいだの力（物がくっつく力、水が水らしい性質をもつ理由）のほとんどは、もとをたどるとクーロンの力ですが、それを計算するには量子力学が要ります（物理シリーズの『原子』（原-7））。
\end{phsaki}

\begin{senmonsho}[電気工学の専門書の書き方]
\begin{bkhonbun}
真空中に置かれた点電荷 $Q_1$, $Q_2$ 間に働く静電力（クーロン力）は
\[ \symbfit{F}=\frac{Q_1Q_2}{4\pi\varepsilon_0r^2}\,\symbfit{a}_r \quad[\mathrm{N}] \]
である．単位正電荷に働く力を電界 $\symbfit{E}$ と定義すれば，点電荷 $Q$ のつくる電界の大きさは $E=Q/(4\pi\varepsilon_0r^2)$ [V/m] となる．
\end{bkhonbun}
\bkyaku{電気工学の本は、電場を「電界」とよび、ベクトルを太字 $\symbfit{F}$・$\symbfit{E}$ で書くことが多い（この本の $\vec F$・$\vec E$）。$\symbfit{a}_r$ は向きを表す長さ1の矢印（この本の $\vec e_r$）。$k$ を使わず、いつも $\dfrac{1}{4\pi\varepsilon_0}$ と書く。「単位正電荷」は、+1\,C の試しの電荷のこと。}
\end{senmonsho}

\begin{taisetsu}
\def\tsl#1{\makebox[7.6\zw][l]{\sffamily\bfseries\small\color{\bkTT}#1}}
\tsl{電荷}電子 $-e$、$e=\qty{1.60e-19}{C}$。電荷は保存する\\[3pt]
\tsl{クーロンの法則}$F=k\,\dfrac{q_1q_2}{r^2}$、$k=\dfrac{1}{4\pi\varepsilon_0}=\qty{8.99e9}{N.m^2/C^2}$\\[3pt]
\tsl{電場}$\vec E=\dfrac{\vec F}{q}$、点電荷 $E=k\,\dfrac{Q}{r^2}$。重ね合わせは矢印の和\\[3pt]
\tsl{双極子}$p=qd$、軸の上の遠くで $E\approx\dfrac{2kp}{r^3}$\\[3pt]
\tsl{広がった電荷}長い線 $\dfrac{\lambda}{2\pi\varepsilon_0r}$、広い面 $\dfrac{\sigma}{2\varepsilon_0}$、正負の2枚のあいだ $\dfrac{\sigma}{\varepsilon_0}$\\[3pt]
\tsl{電場の中}$\vec a=\dfrac{q}{m}\vec E$。一様なら放物線
\end{taisetsu}

\begin{kotoba}
\small
\bkword{帯電}{物が正または負の電気を帯びること。こすると電子が移って起きる。}
\bkword{電荷・電気素量}{電気の量（単位 C）。電子1個の電荷の大きさ $e$ が電気素量。}
\bkword{導体・絶縁体}{電気をよく通す物（金属）と、ほとんど通さない物（ガラス・プラスチック）。}
\bkword{クーロンの法則}{2つの点電荷の力は、電荷の積に比例し、距離の2乗に反比例する。}
\bkword{電場}{そこに置いた電荷が 1\,C あたりに受ける力。単位 N/C（＝V/m）。}
\bkword{電気力線}{電場の向きをつないだ線。こみ合っている所ほど電場が強い。}
\bkword{電気双極子}{$+q$ と $-q$ が少し離れて組になったもの。$p=qd$。}
\bkword{一様な電場}{大きさと向きが、どこでも同じ電場。}
\end{kotoba}

\bkkigocards{\bkkigocard{$Q$・$q$}{キュー}{電荷〔C〕。場をつくる電荷を $Q$、試しの・小さな電荷を $q$。熱量 $Q$（第2巻）とは別の量}{electric charge}{電-1}}
{\bkkigocard{$\vec E$}{イー}{電場〔N/C＝V/m〕。大きさは $E$。ヤング率・エネルギー $E$（第1巻）とは別の量}{electric field}{電-1}}
\bkkigocards{\bkkigocard{$k$}{ケー}{クーロンの定数。約 $\qty{8.99e9}{N.m^2/C^2}$。ばね定数 $k$（第1巻）とは別の量}{Coulomb constant}{電-1}}
{\bkkigocard{$\varepsilon_0$}{イプシロン・ゼロ}{真空の誘電率〔F/m〕。約 $8.85\times10^{-12}$。物質の誘電率 $\varepsilon$ は電-3}{vacuum permittivity}{電-1}}
\bkkigocards{\bkkigocard{$e$}{イー}{電気素量。約 $\qty{1.60e-19}{C}$。電子の電荷は $-e$}{elementary charge}{電-1}}
{\bkkigocard{$\lambda$・$\sigma$}{ラムダ・シグマ}{線電荷密度〔C/m〕・面電荷密度〔C/m$^2$〕。導電率 $\sigma$（電-4）とは別の量}{charge density}{電-1}}
\bkkigocards{\bkkigocard{$\vec p$}{ピー}{双極子モーメント〔C·m〕。大きさ $p=qd$。運動量 $\vec p$（第1巻）とは別の量}{dipole moment}{電-1}}{}

%% ======================== 例題 ========================
\Needspace{10\baselineskip}\bksection{reidai}{例題でたしかめよう}
\label{ph:電-1:ex1}
\begin{phreidai}[1\,C どうしの力]
\qty{1}{C}の点電荷どうしを\qty{1}{m}離して置くと、はたらく力はいくらですか。その力は、何 kg の物にはたらく重力と同じ大きさですか。
\phbamen
2つの点電荷のあいだの力を、クーロンの法則で求める場面です。
\phmushi
電荷をもつ物の大きさ（点電荷とみる）と、周りの空気を無視します。
\phzu
図\ref{ph:電-1:fig-chikara}。$q_1=q_2=\qty{1}{C}$、$r=\qty{1}{m}$。
\phriyu
止まった2つの点電荷の力なので、クーロンの法則（②）。重さとくらべるには、$mg$ と等しいとおく（力-2）。
\phtateru
\[F=k\,\frac{q_1q_2}{r^2},\qquad m=\frac{F}{g}\]
\phtoku
\[F=8.99\times10^9\times\frac{1\times1}{1^2}=\qty{8.99e9}{N}\approx\qty{9.0e9}{N}\]
\[m=\frac{8.99\times10^9}{9.8}\approx\qty{9.2e8}{kg}\]
\phtashikame
単位：N·m²/C²×C²÷m²＝N。\\
大きさ：$9.2\times10^8$\,kg は92万トンで、大きなタンカーの数隻分です。1\,C は、静電気としてはとても大きな電荷です。\\
極端な場合：距離を\qty{1}{km}にしても $\qty{9.0e3}{N}$（約900 kg の重さ）もあります。
\kotae 約 $\qty{9.0e9}{N}$。約 $9.2\times10^8$\,kg（約92万トン）の物の重さと同じ
\end{phreidai}
\begin{jibun}
（1）\qty{1}{\micro C}どうしを\qty{1}{cm}離して置くと、力はいくらですか。（2）同じ形の金属球 A（$+3$\,nC）と B（$-1$\,nC）を触れさせてから離すと、それぞれの電荷はいくらになりますか。
\bkjans{（1）$8.99\times10^9$ $\times(10^{-6})^2\div0.01^2$ で、約\qty{90}{N}。（2）合計 $+2$\,nC を半分ずつ分けて、どちらも $+1$\,nC}
\end{jibun}

\label{ph:電-1:ex2}
\begin{phreidai}[水素原子の中の、電気の力と重力]
水素原子では、陽子の周りを、平均で約 $\qty{5.29e-11}{m}$ 離れて電子が回っています。陽子と電子のあいだの電気の力と、万有引力を求め、くらべましょう。陽子の質量は $\qty{1.67e-27}{kg}$、電子の質量は $\qty{9.11e-31}{kg}$、$G=\qty{6.67e-11}{N.m^2/kg^2}$ とします。
\phbamen
原子の中で、どの力が大切かを見積もる場面です。
\phmushi
電子は陽子から一定の距離にいるとみなし、量子力学の効果（電子の位置のぼやけ）は考えません。
\phzu
陽子 $+e$、電子 $-e$、$r=\qty{5.29e-11}{m}$。
\phriyu
電気の力はクーロンの法則、重力は万有引力の法則（力-2）。どちらも距離の2乗に反比例するので、比は距離によりません。
\phtateru
\[F_{\text{電気}}=k\,\frac{e^2}{r^2},\qquad F_{\text{重力}}=G\,\frac{m_{\mathrm p}m_{\mathrm e}}{r^2}\]
\phtoku
\[F_{\text{電気}}=8.99\times10^9\times\frac{(1.60\times10^{-19})^2}{(5.29\times10^{-11})^2}\approx\qty{8.2e-8}{N}\]
\[F_{\text{重力}}=6.67\times10^{-11}\times\frac{1.67\times10^{-27}\times9.11\times10^{-31}}{(5.29\times10^{-11})^2}\approx\qty{3.6e-47}{N}\]
比は $8.2\times10^{-8}\div(3.6\times10^{-47})\approx2.3\times10^{39}$ です。
\phtashikame
比は $ke^2/(Gm_{\mathrm p}m_{\mathrm e})$ で、$r$ が消えます。どの距離でくらべても同じ比です。\\
大きさ：$10^{39}$ は、1のあとに0が39個並ぶ数です。原子や分子の世界では、重力はまったく無視できます。
\kotae 電気の力 約 $\qty{8.2e-8}{N}$、万有引力 約 $\qty{3.6e-47}{N}$。電気の力は約 $2.3\times10^{39}$ 倍
\end{phreidai}
\begin{jibun}
電子どうしが\qty{1}{nm}（$10^{-9}$\,m）離れているとき、はね返す電気の力と、引き合う万有引力を求めましょう。比はいくらですか。
\bkjans{電気の力 約 $\qty{2.3e-10}{N}$、万有引力 約 $\qty{5.5e-53}{N}$。比は約 $4.2\times10^{42}$}
\end{jibun}

\label{ph:電-1:ex3}
\begin{phreidai}[髪の毛は立ち上がるか（はてな？）]
こすった下じきの電荷を $\qty{-2.0e-8}{C}$、髪の毛1本の先にある電荷を $\qty{+1.0e-10}{C}$ とし、下じきと髪の先が\qty{5}{cm}離れているとします。髪の毛1本は、長さ\qty{10}{cm}、直径\qty{80}{\micro m}の円柱で、密度は\qty{1300}{kg/m^3}とします。髪の先が受ける電気の力と、髪の毛1本の重さをくらべましょう。
\phbamen
はてな？の髪の毛が、下じきに引かれて重力に勝てるかを見積もる場面です。
\phmushi
下じきの電荷は1つの点に集まっているとみなします。ほかの髪の毛からの力と、髪の根元が髪を支える力の向きは考えません。電荷と大きさは、見当をつけるための仮の数です。
\phzu
下じき $Q=\qty{-2.0e-8}{C}$、髪の先 $q=\qty{1.0e-10}{C}$、$r=\qty{0.05}{m}$。髪：半径 $b=\qty{40}{\micro m}$、長さ $l=\qty{0.10}{m}$、密度 $\rho_{\text{髪}}=\qty{1300}{kg/m^3}$。
\phriyu
電気の力はクーロンの法則（②）。重さは、質量（体積×密度）×$g$。
\phtateru
\[F=k\,\frac{|Qq|}{r^2},\qquad m=\rho_{\text{髪}}\times\pi b^2l,\qquad\text{重さ}=mg\]
\phtoku
\[F=8.99\times10^9\times\frac{2.0\times10^{-8}\times1.0\times10^{-10}}{0.05^2}\approx\qty{7.2e-6}{N}\]
\[mg=1300\times\pi\times(40\times10^{-6})^2\times0.10\times9.8\approx\qty{6.4e-6}{N}\]
比は $7.2\div6.4\approx1.1$ です。
\phtashikame
単位：どちらも N。\\
大きさ：髪の毛1本の質量は約 $\qty{6.5e-7}{kg}$（0.65 mg）で、とても軽いのです。髪の先の電荷 $10^{-10}$\,C は、電子にして約6億個分にすぎませんが、それで重さと同じくらいの力が出ます。\\
極端な場合：下じきを遠ざけると、力は $1/r^2$ で弱まります（じぶんでとこう）。
\kotae 電気の力 約 $\qty{7.2e-6}{N}$、重さ 約 $\qty{6.4e-6}{N}$。電気の力が約1.1倍で、髪の毛は持ち上がる
\end{phreidai}
\begin{jibun}
例題3で、下じきを\qty{10}{cm}まで遠ざけると、電気の力はいくらになりますか。重さの何倍ですか。電気の力が重さと同じになる距離も求めましょう。
\bkjans{距離2倍で $\tfrac14$ の約 $\qty{1.8e-6}{N}$。重さの約0.28倍で、髪は下がる。同じになるのは約\qty{5.3}{cm}}
\end{jibun}

\label{ph:電-1:ex4}
\begin{phreidai}[2つの電荷の電場を矢印でたす]
$+\qty{1}{\micro C}$ を原点 $(0,\ 0)$ に、$-\qty{1}{\micro C}$ を $(0.1,\ 0)$〔m〕に置きます。点 $(0.05,\ 0.05)$ の電場の大きさと向きを求めましょう。
\phbamen
2つの点電荷の電場を、重ね合わせで求める場面です（図\ref{ph:電-1:fig-kasane}と同じ形）。
\phmushi
周りの物の影響を無視し、点電荷とみます。
\phzu
点 P$(0.05,\ 0.05)$ は、どちらの電荷からも $\sqrt{0.05^2+0.05^2}\approx\qty{0.0707}{m}$ で、2つの電荷を結ぶ線から45°の向きにあります。
\phriyu
電場は、それぞれの点電荷の電場の矢印の和（③④）。対称な置き方なので、成分に分けると打ち消し合う部分があります。
\phtateru
1つの電荷の電場の大きさ $E_1=k\,|Q|/r^2$。$+$ の電場は原点から P へ向かう向き $(\cos45^\circ,\ \sin45^\circ)$、$-$ の電場は P から $(0.1,\ 0)$ へ向かう向き $(\cos45^\circ,\ -\sin45^\circ)$。
\phtoku
\[E_1=8.99\times10^9\times\frac{1\times10^{-6}}{0.0707^2}=8.99\times10^9\times\frac{10^{-6}}{0.005}\approx\qty{1.80e6}{N/C}\]
\[E_x=2E_1\cos45^\circ=2\times1.80\times10^6\times0.707\approx\qty{2.5e6}{N/C}\]
\[E_y=E_1\sin45^\circ-E_1\sin45^\circ=0\]
\phtashikame
向き：たての成分が打ち消し合い、$+$ から $-$ へ向かう向き（$x$ の正の向き）。電気力線の図\ref{ph:電-1:fig-rikisen}（左）で、正と負のあいだの上の線が右向きなのと合います。\\
大きさ：1つの電場の $\sqrt2\approx1.41$ 倍です（45°の2本の矢印の和）。
\kotae 大きさ 約 $\qty{2.5e6}{N/C}$、向きは $x$ の正の向き（$+$ の電荷から $-$ の電荷へ平行に）
\end{phreidai}
\begin{jibun}
例題4で、$(0.1,\ 0)$ の電荷も $+\qty{1}{\micro C}$ にすると、点 $(0.05,\ 0.05)$ の電場の大きさと向きはどうなりますか。
\bkjans{横の成分が打ち消し合い、$y$ の正の向き（真上）に約 $\qty{2.5e6}{N/C}$}
\end{jibun}

\label{ph:電-1:ex5}
\begin{phreidai}[力が0になる点]
$+4q$ の点電荷を $x=0$ に、$+q$ の点電荷を $x=\qty{0.3}{m}$ に置きます（$q>0$）。2つを結ぶ線の上で、ほかの電荷を置いても力を受けない点（電場が0の点）はどこですか。
\phbamen
2つの電荷の電場が打ち消し合う点を探す、高校の入試にもよく出る形です。
\phmushi
点電荷とみます。
\phzu
求める点を $x$ とします。2つの電荷のあいだ（$0<x<0.3$）では、$+4q$ の電場は右向き、$+q$ の電場は左向きで、打ち消し合えます。外側では、どちらも同じ向きになるので、打ち消し合いません。
\phriyu
重ね合わせ（③）。大きさが等しく向きが反対なら、電場は0。
\phtateru
\[k\,\frac{4q}{x^2}=k\,\frac{q}{(0.3-x)^2}\]
\phtoku
両辺の $kq$ を消して、平方根をとります（どちらも正の数なので）。
\[\frac{2}{x}=\frac{1}{0.3-x}\qquad\text{より}\qquad2(0.3-x)=x,\qquad x=\qty{0.2}{m}\]
\phtashikame
$x=0.2$ では、$4q/0.2^2=100q$、$q/0.1^2=100q$ で、等しくなっています。\\
大きいほうの電荷から遠い所（小さいほうの電荷の近く）にあるのは、もっともらしい答えです。
\kotae $+4q$ から\qty{0.2}{m}（$+q$ から\qty{0.1}{m}）の点
\end{phreidai}
\begin{jibun}
$+q$ を $x=0$ に、$+9q$ を $x=\qty{0.4}{m}$ に置きます。2つのあいだで電場が0の点を求めましょう。
\bkjans{$\dfrac1x=\dfrac{3}{0.4-x}$ より $x=\qty{0.1}{m}$（$+q$ から\qty{0.1}{m}）}
\end{jibun}

\label{ph:電-1:ex6}
\begin{phreidai}[正三角形の頂点の電荷]
1辺\qty{0.1}{m}の正三角形の3つの頂点に、$+\qty{1}{\micro C}$ ずつの点電荷を置きました。1つの電荷が受ける力の大きさと向きを求めましょう。
\phbamen
3つ以上の電荷の力を、重ね合わせで求める場面です。
\phmushi
点電荷とみます。
\phzu
1つの頂点 A は、ほかの2つ B・C から、それぞれ辺にそって外向きの力を受けます。2つの力のあいだの角は60°です。
\phriyu
1組の力はクーロンの法則、合力は矢印の和（③）。2つの力は大きさが同じなので、合力は角を二等分する向き（三角形の中心から A へ向かう向き）。
\phtateru
\[F_1=k\,\frac{q^2}{a^2},\qquad F=2F_1\cos30^\circ\]
\phtoku
\[F_1=8.99\times10^9\times\frac{(1\times10^{-6})^2}{0.1^2}\approx\qty{0.90}{N},\qquad F=2\times0.899\times0.866\approx\qty{1.6}{N}\]
\phtashikame
2つの力のなす角60°を二等分すると、それぞれ30°なので、$\cos30^\circ$ をかけた成分だけが残ります。合力は1組の力の $\sqrt3\approx1.73$ 倍です。
\kotae 約\qty{1.6}{N}、三角形の中心から外向き
\end{phreidai}
\begin{jibun}
1辺\qty{0.1}{m}の正方形の4つの頂点に、$+\qty{1}{\micro C}$ ずつ置きました。1つの電荷が受ける力の大きさを求めましょう。
\bkjans{となり2つから\qty{0.90}{N}ずつ（直角）、対角から\qty{0.45}{N}。合わせて約\qty{1.7}{N}、中心から外向き}
\end{jibun}

\label{ph:電-1:ex7}
\begin{phreidai}[点電荷の電場の表]
\qty{1}{nC}の点電荷から、\qty{0.1}{m}・\qty{0.2}{m}・\qty{0.3}{m}・\qty{0.5}{m}・\qty{1}{m}の所の電場の大きさを求め、表にしましょう。距離と電場の関係を表で確かめましょう。
\phbamen
1つの点電荷の周りの電場の強さを、場所ごとに求める場面です（8-5⑥ と同じ数）。
\phmushi
点電荷とみます。
\phzu
$Q=\qty{1e-9}{C}$。
\phriyu
点電荷の電場 $E=kQ/r^2$（④）。
\phtateru
\[E=k\,\frac{Q}{r^2}=\frac{8.99}{r^2}\quad\text{（$kQ=8.99\times10^9\times10^{-9}=8.99$）}\]
\phtoku
$r=0.1$：$8.99\div0.01=899$、$r=0.2$：$8.99\div0.04\approx225$、$r=0.3$：$8.99\div0.09\approx99.9$、$r=0.5$：$8.99\div0.25\approx36.0$、$r=1$：8.99（単位は N/C）。表\ref{ph:電-1:tab-1nC}のとおりです。
\phtashikame
0.1 と 0.2：比は $899\div225\approx4$。0.1 と 0.3：$899\div99.9\approx9$。距離が2倍で $\tfrac14$、3倍で $\tfrac19$ になっています。
\kotae 899、225、99.9、36.0、8.99（N/C）
\end{phreidai}
\begin{jibun}
（1）\qty{2}{nC}の点電荷から\qty{0.1}{m}・\qty{0.3}{m}の所の電場を求めましょう。（2）\qty{1}{nC}の点電荷の電場が\qty{100}{N/C}になるのは、何 m の所ですか。
\bkjans{（1）約\qty{1.8e3}{N/C}、約\qty{200}{N/C}。（2）$r=\sqrt{8.99\div100}$ で、約\qty{0.30}{m}}
\end{jibun}

\label{ph:電-1:ex8}
\begin{phreidai}[双極子の電場は、距離の3乗で弱まる]
$+\qty{1}{nC}$ と $-\qty{1}{nC}$ を\qty{1}{cm}離した双極子があります。軸の上で、まん中から\qty{0.1}{m}と\qty{0.2}{m}の所（$+$ の側）の電場を、重ね合わせの式と近似の式の両方で求め、2つの距離の電場の比を求めましょう。
\phbamen
遠くから見た双極子の電場の弱まり方を確かめる場面です。
\phmushi
点電荷とみます。
\phzu
図\ref{ph:電-1:fig-soukyoku}。$q=\qty{1e-9}{C}$、$d=\qty{0.01}{m}$、$p=qd=\qty{1e-11}{C.m}$。
\phriyu
正確には2つの点電荷の重ね合わせ（③）。$r\gg d$ では $E\approx2kp/r^3$（⑥）。
\phtateru
\[E=kq\Bigl(\frac{1}{(r-d/2)^2}-\frac{1}{(r+d/2)^2}\Bigr),\qquad E\approx\frac{2kp}{r^3}\]
\phtoku
$r=0.1$ のとき（$kq=8.99$、$2kp=0.1798$）：
\[E=8.99\Bigl(\frac{1}{0.095^2}-\frac{1}{0.105^2}\Bigr)\approx996.1-815.4\approx\qty{181}{N/C}\]
近似の式では $\dfrac{0.1798}{0.1^3}\approx\qty{180}{N/C}$。\\
$r=0.2$ のとき：
\[E=8.99\Bigl(\frac{1}{0.195^2}-\frac{1}{0.205^2}\Bigr)\approx\qty{22.5}{N/C}\]
近似の式でも $\dfrac{0.1798}{0.2^3}\approx\qty{22.5}{N/C}$。\\
比は $181\div22.5\approx8.0$ です。
\phtashikame
距離2倍で $\tfrac18$ なので、3乗に反比例しています。1つの電荷だけの電場は、\qty{0.1}{m}で899\,N/C（例題7）なので、双極子の電場は、その約5分の1しかありません。
\kotae \qty{0.1}{m}で約\qty{181}{N/C}（近似180）、\qty{0.2}{m}で約\qty{22.5}{N/C}。比は約8
\end{phreidai}
\begin{jibun}
$\pm\qty{2}{nC}$ を\qty{1}{cm}離した双極子の軸の上で、まん中から\qty{0.1}{m}と\qty{0.3}{m}の所の電場を、近似の式で求めましょう。
\bkjans{\qty{0.1}{m}：約\qty{360}{N/C}、\qty{0.3}{m}：約\qty{13}{N/C}（27分の1）}
\end{jibun}

\label{ph:電-1:ex9}
\begin{phreidai}[長い線の電荷の電場]
まっすぐな長い線に、1 m あたり\qty{1}{nC}の電荷が一様に広がっています。線から\qty{0.1}{m}の所の電場を求めましょう。
\phbamen
帯電した長い電線の周りの電場を求める場面です。
\phmushi
線は、\qty{0.1}{m}よりずっと長い（限りなく長い）とします。線の太さを無視します。
\phzu
図\ref{ph:電-1:fig-sen}。$\lambda=\qty{1e-9}{C/m}$、$r=\qty{0.1}{m}$。
\phriyu
小さく分けた電荷の電場を積分した結果 $E=2k\lambda/r$（⑦）を使います。
\phtateru
\[E=\frac{2k\lambda}{r}\]
\phtoku
\[E=\frac{2\times8.99\times10^9\times1\times10^{-9}}{0.1}\approx\qty{180}{N/C}\]
\phtashikame
$\dfrac{\lambda}{2\pi\varepsilon_0r}=\dfrac{10^{-9}}{2\pi\times8.85\times10^{-12}\times0.1}\approx180$ でも同じです。\\
大きさ：同じ\qty{0.1}{m}でも、\qty{1}{nC}の点電荷（899）より弱いのは、線の電荷の多くが遠くにあるからです。距離を2倍にすると半分（$1/r$）になります。電-2 例題2 では、ガウスの法則で、同じ答えを積分なしに出します。
\kotae 約\qty{180}{N/C}（線から外向き）
\end{phreidai}
\begin{jibun}
1 m あたり\qty{2}{nC}の長い線電荷から、\qty{0.05}{m}と\qty{0.2}{m}の所の電場を求めましょう。
\bkjans{\qty{0.05}{m}：約\qty{720}{N/C}、\qty{0.2}{m}：約\qty{180}{N/C}（距離4倍で $\tfrac14$）}
\end{jibun}

\label{ph:電-1:ex10}
\begin{phreidai}[円環の軸の上で、いちばん強い所]
半径\qty{0.1}{m}の細い輪に、$\qty{1e-8}{C}$ の電荷が一様に広がっています。軸の上で電場がいちばん強くなる位置と、そのときの電場を求めましょう。
\phbamen
広がった電荷の電場を積分で求め、微分で最大の点を探す場面です。
\phmushi
輪の太さを無視します。
\phzu
図\ref{ph:電-1:fig-enkan}。$Q=\qty{1e-8}{C}$、$a=\qty{0.1}{m}$。
\phriyu
円環の軸の上の電場（⑦）を $z$ で微分して0になる所が、いちばん強い所（5-7）。
\phtateru
\[E=\frac{kQz}{(z^2+a^2)^{3/2}},\qquad\frac{\dd E}{\dd z}=0\ \text{より}\ z=\frac{a}{\sqrt2}\]
\phtoku
\[z=\frac{0.1}{\sqrt2}\approx\qty{0.0707}{m},\qquad z^2+a^2=0.005+0.01=0.015\]
\[E=\frac{8.99\times10^9\times10^{-8}\times0.0707}{0.015^{3/2}}=\frac{6.36}{1.84\times10^{-3}}\approx\qty{3.5e3}{N/C}\]
\phtashikame
$z=0$ で0、遠くで $kQ/z^2$ に近づくので、途中に最大があるのは当然です。$\dd E/\dd z$ の分子 $a^2-2z^2$ は、$z<a/\sqrt2$ で正（増える）、それより先で負（減る）です。
\kotae 中心から約\qty{0.071}{m}、約\qty{3.5e3}{N/C}
\end{phreidai}
\begin{jibun}
例題10の円環で、軸の上の\qty{1}{m}の所の電場を、円環の式と、点電荷とみた式で求め、くらべましょう。
\bkjans{円環の式 約\qty{88.6}{N/C}、点電荷とみると\qty{89.9}{N/C}（約1.5\,\% の違い）}
\end{jibun}

\label{ph:電-1:ex11}
\begin{phreidai}[広い平らな面の電場]
広い平らな面に、1 m$^2$ あたり\qty{1}{\micro C}の電荷が一様に広がっています。面から\qty{1}{cm}・\qty{5}{cm}の所の電場を求めましょう。
\phbamen
帯電した広い板の近くの電場を求める場面です。
\phmushi
面はどこまでも広いとみなします（板の大きさより、ずっと近い所を考える）。
\phzu
図\ref{ph:電-1:fig-2mai}の、正の面1枚だけ。$\sigma=\qty{1e-6}{C/m^2}$。
\phriyu
広い平らな面の電場 $\sigma/(2\varepsilon_0)$ は、距離によりません（⑦）。
\phtateru
\[E=\frac{\sigma}{2\varepsilon_0}\]
\phtoku
\[E=\frac{1\times10^{-6}}{2\times8.85\times10^{-12}}\approx\qty{5.6e4}{N/C}\]
\phtashikame
\qty{1}{cm}でも\qty{5}{cm}でも同じです。ただし、本物の板（たとえば半径\qty{0.1}{m}の円板）では、\qty{5}{cm}離れると広い面の値の約55\,\% に下がります（⑦の2倍にすると？）。
\kotae どちらも約 $\qty{5.6e4}{N/C}$（面に垂直に外向き）
\end{phreidai}
\begin{jibun}
1 m$^2$ あたり\qty{5}{\micro C}の広い面1枚の電場を求めましょう。また、$+\qty{5}{\micro C/m^2}$ と $-\qty{5}{\micro C/m^2}$ の広い面を向かい合わせたとき、あいだと外側の電場を求めましょう。
\bkjans{1枚：約 $\qty{2.8e5}{N/C}$。向かい合わせ：あいだ約 $\qty{5.6e5}{N/C}$、外側0}
\end{jibun}

\label{ph:電-1:ex12}
\begin{phreidai}[インクジェットの滴を曲げる]
質量 $\qty{1e-10}{kg}$、電荷 $\qty{2e-13}{C}$ のインクの滴が、速さ\qty{20}{m/s}で、長さ\qty{1}{cm}の2枚の板のあいだに、電場に垂直に入ります。板のあいだの電場は $\qty{1e6}{N/C}$ です。板を出るまでに、滴は横に何 mm ずれますか。
\phbamen
一様な電場の中を通る、帯電した粒の動きを求める場面です。
\phmushi
重力と空気の抵抗、板のはしで電場が弱まる所を無視します。数はプリンターの見当をつけるための仮の数です。
\phzu
図\ref{ph:電-1:fig-ink}。$m=\qty{1e-10}{kg}$、$q=\qty{2e-13}{C}$、$E=\qty{1e6}{N/C}$、$l=\qty{0.01}{m}$、$v_0=\qty{20}{m/s}$。
\phriyu
一様な電場では加速度 $qE/m$ が一定で、力-1 の放物運動と同じ（⑧）。
\phtateru
\[a=\frac{qE}{m},\qquad t=\frac{l}{v_0},\qquad y=\frac12at^2\]
\phtoku
\[a=\frac{2\times10^{-13}\times10^6}{10^{-10}}=\qty{2000}{m/s^2},\qquad t=\frac{0.01}{20}=\qty{5e-4}{s}\]
\[y=\frac12\times2000\times(5\times10^{-4})^2=\qty{2.5e-4}{m}=\qty{0.25}{mm}\]
\phtashikame
単位：C×N/C÷kg＝N/kg＝m/s²。\\
重力との比：$2000\div9.8\approx200$ 倍なので、重力を無視してよかったのです。\\
板を出るときの横の速さは $at=\qty{1}{m/s}$ で、そのあとまっすぐ進むので、紙までの距離が長いほど、紙の上でのずれは大きくなります。
\kotae 約\qty{0.25}{mm}
\end{phreidai}
\begin{jibun}
電子が、速さ $\qty{1e7}{m/s}$ で、長さ\qty{2}{cm}の板のあいだに、電場 $\qty{1e4}{N/C}$ に垂直に入ります。板を出るまでのずれを求めましょう（電子の質量 $\qty{9.11e-31}{kg}$）。
\bkjans{$a=eE/m\approx\qty{1.76e15}{m/s^2}$、$t=\qty{2e-9}{s}$、ずれ 約\qty{3.5}{mm}}
\end{jibun}

%% ======================== まちがえやすいところ ========================
\Needspace{11\baselineskip}\bksection{tsumazuki}{まちがえやすいところをチェック}
\begin{tsumazuki}{距離が2倍なら、力は半分？}
\bkng{クーロンの力は距離に反比例するので、距離2倍で半分。}
\bkok{距離の\textbf{2乗}に反比例するので、距離2倍で $\tfrac14$。}
長い線の電場（$1/r$）、双極子の電場（$1/r^3$）、広い面の電場（距離によらない）と、電荷の形で弱まり方が違います。点電荷は $1/r^2$ です（②⑥⑦）。
\end{tsumazuki}

\begin{tsumazuki}{電場は、力を受ける電荷で変わる？}
\bkng{試しの電荷を2倍にすると、受ける力が2倍になるので、電場も2倍になる。}
\bkok{力も2倍、$q$ も2倍なので、$E=F/q$ は変わらない。電場は、周りの電荷と場所で決まる。}
電場は「1\,C あたりの力」なので、試しの電荷によりません（④）。
\end{tsumazuki}

\begin{tsumazuki}{電場を大きさだけでたす}
\bkng{例題4で、2つの電場の大きさ $1.8\times10^6$ をたして $3.6\times10^6$\,N/C。}
\bkok{2つの電場は向きが違うので、矢印としてたす。たての成分が打ち消し合って、$2.5\times10^6$\,N/C になる。}
電場も力も向きをもつ量です。成分に分けてたします（③）。
\end{tsumazuki}

\begin{tsumazuki}{負の電荷も、電場の向きに動く？}
\bkng{電子は、電場の向きに力を受ける。}
\bkok{$\vec F=q\vec E$ で $q<0$ なので、電子は電場と\textbf{反対}の向きに力を受ける。}
電場の向きは「正の電荷が受ける力の向き」と決めてあります（④）。
\end{tsumazuki}

\begin{tsumazuki}{電気力線は、電荷が動く道すじ？}
\bkng{止まっていた電荷を放すと、電気力線にそって動く。}
\bkok{電気力線の向きは、その点の\textbf{加速度}の向き。動き出した電荷は勢いがつくので、曲がった力線からははずれていく。}
力線にそって動くのは、まっすぐな力線の上か、とても遅く動く（水の中のちりのように、すぐに止められる）ときだけです（⑤⑧）。
\end{tsumazuki}

\begin{column}{1\,C は、どれくらいの電気？}
\qty{1}{C}は、\qty{1}{A}の電流が1秒間に運ぶ電荷です。スマホの電池に「\qty{3000}{mAh}」と書いてあれば、\qty{3}{A}（\qty{3000}{mA}）の電流を1時間（\qty{3600}{s}）流せる電荷、つまり $3\times3600=\num{10800}$\,C を送り出せる、という意味です（電池の中の電荷が全部なくなるのではなく、電池の中の化学反応で、電荷をくみ上げて回路に送り続けます：電-4）。いっぽう、静電気で物にたまる電荷は、\qty{1}{nC}から\qty{1}{\micro C}くらいです。同じ「クーロン」でも、静電気と電流では、桁が1兆倍以上も違います。それでも、静電気の力が感じられるのは、たまった電荷が正だけ・負だけにかたよっているからです。電流では、導線の中の電子と、同じ数の原子核の正の電荷がいっしょにあり、外から見ると打ち消し合っています。
\end{column}

%% ======================== 工学ではこう使う ========================
\Needspace{18\baselineskip}\bksection{phgear}{工学ではこう使う}
\begin{phkougaku}[静電気で、物を集める・動かす]
\textbf{電気集じん器}：発電所や工場の煙の中の細かい粒を取り除く装置です。細い電線（放電線）に高い電圧をかけて、その周りの空気に電荷をつくり、通りすぎる粒を帯電させます。帯電した粒は、反対の符号の大きな板（集じん板）へ電場で引かれて、くっつきます。粒が小さいほど重力や空気の流れの影響を受けにくくなるので、電場で集める方法が役に立ちます。

\textbf{静電塗装}：塗料の霧に電荷をのせると、霧の粒どうしがはね返し合って均一に広がり、反対の符号にした車体の裏側にも回りこみます。むだになる塗料を減らせます。

\textbf{複写機・レーザープリンター}：光で電荷のもようを作り、帯電した粉（トナー）をそのもように合わせてつけます。粉の粒はとても小さく軽いので、わずかな電荷でも、電気の力が重さよりずっと大きくなります（例題3・12と同じ考え）。

\begin{phketa}[静電気の数の目安]
\phketaline{静電気でたまる電荷}{\qty{1}{nC}〜\qty{1}{\micro C}（ストロー・下じき）}
\phketaline{パチッと感じる体の電圧}{約\qty{3}{kV}以上}
\phketaline{雷1回で運ばれる電荷}{数 C〜数十 C}
\phketaline{1\,nC から\qty{0.1}{m}の電場}{\qty{899}{N/C}}
\phboxhead{答えが桁で合っているか}
$k=9\times10^9$ はとても大きいので、1\,C どうしの力はとんでもなく大きくなります。静電気の問題で、力が何千 N も出てきたら、電荷の単位（\unit{\micro C}・nC の $10^{-6}$・$10^{-9}$）を入れ忘れていないか確かめましょう。
\end{phketa}
\end{phkougaku}
\phstation{高電圧工学,静電気工学,電子写真,環境装置}{帯電、クーロン力、電界、電気力線、静電誘導、静電気放電（ESD）、電気集じん}

%% ======================== 章末問題 ========================
\bkendproblems
\bkproblemgroup{基本}{3}
\bkproblem 同じ形の小さな金属球 A（$+\qty{3.0}{nC}$）と B（$-\qty{2.0}{nC}$）を、\qty{3.0}{cm}離して置きました。（1）2つのあいだの力の大きさを求め、引き合うかはね返すか答えましょう。（2）A と B を触れさせてから、もとの\qty{3.0}{cm}にもどしました。力の大きさと向きはどうなりますか。
\bkpans{（1）$kq_1q_2/r^2=8.99\times10^9$ $\times6.0\times10^{-18}$ $\div0.030^2$ で、約 $\qty{6.0e-5}{N}$。引き合う。（2）どちらも $+\qty{0.5}{nC}$ になり、約 $\qty{2.5e-6}{N}$ ではね返す}
\bkproblem $+\qty{4}{nC}$ の点電荷から\qty{0.2}{m}の所の電場の大きさを求めましょう。そこに電子を置くと、どちら向きに、どれだけの力を受けますか。
\bkpans{約\qty{899}{N/C}（外向き）。電子は電荷へ向かう向きに、約 $\qty{1.4e-16}{N}$}
\bkproblem $+q$ を $x=0$ に、$+4q$ を $x=\qty{0.6}{m}$ に置きます。2つのあいだで、電場が0になる点を求めましょう。
\bkpans{$\dfrac{1}{x}=\dfrac{2}{0.6-x}$ より $x=\qty{0.2}{m}$（$+q$ から\qty{0.2}{m}）}
\bkproblem $+\qty{2}{\micro C}$ の点電荷を $(0,\ 0)$ と $(0.2,\ 0)$〔m〕に置きます。点 $(0.1,\ 0.1)$ の電場の大きさと向きを求めましょう。
\bkpans{1つの電場は約 $\qty{9.0e5}{N/C}$。横は打ち消し、真上（$y$ の正の向き）に約 $\qty{1.3e6}{N/C}$}
\bkproblem 電子が速さ $\qty{2.0e6}{m/s}$ で、長さ\qty{5.0}{cm}の2枚の板のあいだに、電場 \qty{200}{N/C} に垂直に入ります。（1）電子の加速度、（2）板を出るまでのずれを求めましょう。（3）電子にはたらく重力を無視してよいわけを、加速度と $g$ をくらべて説明しましょう。
\bkpans{（1）$eE/m=1.60\times10^{-19}\times200$ $\div(9.11\times10^{-31})$ で、約 $\qty{3.5e13}{m/s^2}$。（2）$t=\qty{2.5e-8}{s}$、ずれ約\qty{1.1}{cm}。（3）$g$ の約 $3.6\times10^{12}$ 倍で、重力の効果はとても小さい}
\bkproblem 原子核の中では、陽子どうしが約 $\qty{1e-15}{m}$ まで近づいています。このときの陽子どうしのクーロンの力を求めましょう。原子核がばらばらにならないのはなぜだと考えられますか。
\bkpans{$8.99\times10^9$ $\times(1.60\times10^{-19})^2$ $\div(10^{-15})^2$ で、約\qty{230}{N}のはね返す力。これより強い別の引き合う力（核力）がはたらいているから（原子核の力は物理シリーズの『原子』で学ぶ）}
\bkproblemgroup{チャレンジ}{4}
\bkproblem $\pm\qty{1}{nC}$ を\qty{1}{cm}離した双極子の、垂直二等分線の上で、まん中から\qty{0.1}{m}の所の電場を、重ね合わせで正確に求めましょう。また、$r\gg d$ のときの近似の式 $E\approx kp/r^3$ を導き、軸の上（例題8）の電場とくらべましょう（例題4と例題8を組み合わせる）。
\bkpans{2つの電場のたての成分は打ち消し、軸に平行な成分が残る。$E=\dfrac{kqd}{(r^2+d^2/4)^{3/2}}$ で、約\qty{89.6}{N/C}。$r\gg d$ で $kqd/r^3=kp/r^3\approx\qty{89.9}{N/C}$。軸の上の約半分で、向きは $\vec p$ と反対}
\bkproblem 長さ\qty{1.0}{m}のまっすぐな線に、1 m あたり\qty{1}{nC}の電荷が一様にあります。線のまん中から、線に垂直に\qty{0.1}{m}の所の電場を、⑦の式で求めましょう。限りなく長い線の値（例題9）の何\,\% ですか。
\bkpans{$\dfrac{2k\lambda}{r}\times\dfrac{0.5}{\sqrt{0.1^2+0.5^2}}\approx180\times0.981$ で、約\qty{176}{N/C}。約98\,\%}
\bkproblem 半径\qty{0.1}{m}の円板に、1 m$^2$ あたり\qty{1}{\micro C}の電荷が一様にあります。軸の上で円板から\qty{1}{cm}の所の電場を、⑦の円板の式で求めましょう。広い面の値（例題11）の何\,\% ですか。
\bkpans{$5.65\times10^4\times\Bigl(1-\dfrac{0.01}{\sqrt{0.01^2+0.1^2}}\Bigr)$ で、約 $\qty{5.1e4}{N/C}$。約90\,\%}
\begin{furikaeri}
\begin{checklist}
\item こすると物が電気を帯びるわけを、電子が移ることで説明できる。\bkfurref{ph:電-1:sec1}
\item クーロンの法則と重ね合わせで、いくつかの電荷のあいだの力と、電場を求められる。\bkfurref{ph:電-1:sec3}
\item 一様な電場の中の電荷の動きを、力学の放物運動と同じ考えで求められる。\bkfurref{ph:電-1:sec8}
\item 線・円環・平面に広がった電荷の電場を積分で求め、電気双極子の電場が距離の3乗に反比例して弱まることを示せる。\bkfurref{ph:電-1:sec7}
\end{checklist}
\end{furikaeri}
\bkfukidashi{mina}{髪の毛の先の、ほんの少しの電荷で、髪の重さと同じくらいの力が出るんだね。電気の力って、重力よりずっと強いんだ！}
